\documentclass[11pt]{article}
\usepackage[a4paper,margin=1in]{geometry}
\usepackage{fontspec}
\setmainfont{Times New Roman}
\usepackage{amsmath,amssymb,mathtools}
\usepackage{booktabs,longtable,tabularx,array,multirow,threeparttable,threeparttablex,makecell}
\usepackage{graphicx,adjustbox,pdflscape}
\usepackage{algorithm}
\usepackage{algpseudocode}
\usepackage{enumitem}
\usepackage{xurl}
\usepackage[colorlinks=true,allcolors=blue]{hyperref}
\usepackage{caption}
\newcolumntype{Y}{>{\raggedright\arraybackslash}X}
\numberwithin{table}{section}
\numberwithin{figure}{section}
\setlength{\parindent}{1.5em}
\setlength{\parskip}{0.35em}
\sloppy
\title{A user dissatisfaction-aware cluster-first route-second framework for the static bike repositioning problem}
\author{}
\date{}
\begin{document}
\maketitle

\begin{abstract}
This paper addresses the multi-vehicle static bike repositioning problem (BRP) with the user dissatisfaction function as the service measure. A cluster-first route-second framework is proposed in which the clustering phase is formulated as a mixed-integer linear program that jointly maximizes the estimated total user dissatisfaction reduction and enforces per-vehicle routing time feasibility through a maximum spanning star approximation of within-cluster travel. Each cluster is then solved as an independent single-vehicle BRP on its assigned stations. Because the clustering model is intractable at city scale, a Hybrid Genetic Search for UDF-aware Clustering (HGS-UC) is developed to solve it. Within HGS-UC, each candidate clustering is evaluated by an activated transfer selection (ATS) procedure, which estimates a time-feasible set of surplus-to-deficit bike transfers for the cluster by selecting transfers in decreasing order of user dissatisfaction reduction per unit handling time. Computational experiments on instances generated from New York Citi Bike data show that, on networks of 50 and 100 stations, HGS-UC obtains near-optimal solutions within seconds to a few minutes; on most of the 100-station instances, its mean solution exceeds the best feasible solution returned by a commercial solver applied directly to the arc-indexed formulation under a two-hour limit. On instances of 200 to 1,000 stations, HGS-UC achieves statistically significant improvements over five baseline clustering methods across all tested combinations of network size, repositioning time budget, and vehicle capacity, with average UDF reduction gains of 14.73\% to 40.67\% under the default time budget and vehicle capacity, and consistent improvements across variations in both parameters.
\end{abstract}

\section{Introduction}

Bike-sharing systems (BSSs) have been widely adopted as a sustainable urban mobility option, providing short-distance trips and first- and last-mile connectivity to public transit (Shui and Szeto, 2020). A persistent operational challenge in these systems is the spatial and temporal mismatch between bike supply and dock availability at individual stations. During peak commuting hours, asymmetric travel patterns leave some stations nearly full and others nearly empty, so that users arriving to return or rent a bike often find no available dock or bike (Hu et al., 2025). Persistent failures in providing sufficient bikes or docks reduce user confidence in BSSs and may ultimately lead to abandonment of the system (Raviv et al., 2013).

To mitigate this imbalance, operators deploy dedicated vehicles to redistribute bikes among stations. The problem of routing these vehicles and determining the pick-up and drop-off quantities at each station is known as the bike repositioning problem (BRP) (Ho and Szeto, 2014). Depending on when and how the operation is carried out, the BRP is classified into static and dynamic variants (Raviv et al., 2013; Shui and Szeto, 2020). In the static variant, repositioning tours are executed during an off-peak period (often overnight) under the assumption that user activity is negligible, with the goal of rearranging station inventories in preparation for the following day (Raviv et al., 2013). By contrast, dynamic repositioning is performed during service hours, with decisions updated periodically in response to time-varying demand and evolving station states (Zhang et al., 2017). In this paper, we focus on the static BRP.

The objectives considered in static BRP formulations generally involve minimizing operating cost, demand-side service measures, or a weighted combination of both (Shui and Szeto, 2020). Operating cost is typically expressed as total vehicle travel time or distance (Benchimol et al., 2011; Chemla et al., 2013). Some formulations also include bike loading and unloading times (Raviv et al., 2013; Szeto et al., 2016). Demand-side service is commonly quantified by specifying a target inventory level or interval for each station and penalizing deviations of the inventory from that reference after repositioning. Such penalties take various forms: total absolute deviation from target fill levels (Di Gaspero et al., 2013; Rainer-Harbach et al., 2015), total unmet demand (Szeto et al., 2016), or a quadratic function of the gap between the resulting and ideal inventory levels (Ho and Szeto, 2014).

However, minimizing deviation from a target level does not necessarily minimize the expected number of service failures. Raviv and Kolka (2013) modeled each station in BSSs as a finite-capacity queueing system with non-homogeneous Poisson rental and return arrivals. They defined the User Dissatisfaction Function (UDF) as the expected number of users who fail to rent or return a bike over the planning horizon and expressed it as a function of the station’s inventory after repositioning. They further proved the convexity of UDF. As a result, two repositioning plans that achieve similar deviations from target inventories can yield different expected numbers of rental and return failures. Therefore, the UDF provides a more accurate basis for evaluating repositioning outcomes than deviation-based penalties.

Studies that incorporate the UDF in the objective remain scarce compared with those that adopt deviation-based measures. Raviv et al. (2013) were the first to present MILP formulations for the static BRP with UDF in the objective. Their formulations minimized a weighted sum of total UDF and vehicle travel time and were solved with CPLEX using a two-phase acceleration strategy. Forma et al. (2015) minimized a weighted sum of UDF and travel distance and solved the problem with a three-step math-heuristic that clusters stations, routes vehicles through clusters with preliminary (un)loading decisions, and finally re-solves the original repositioning problem with the restriction that vehicle traversal is limited to within a cluster or between consecutive clusters. Ho and Szeto (2017) adopted UDF as the penalty function in their formulation and solved the model using a hybrid large neighbourhood search (H-LNS). In the dynamic setting, Zhang et al. (2017) built a time-space network flow model that minimizes the weighted sum of vehicle travel cost and UDF. They developed a heuristic that combines LP relaxation with set-covering-based route selection and a dynamic programming subroutine for loading sequence optimization and extended the framework through a rolling horizon scheme for multiple repositioning periods.

Solving the BRP on large-scale networks remains computationally challenging due to its NP-hardness (Benchimol et al., 2011). Exact methods such as branch-and-cut have been applied to the problem, but they become intractable for large and realistic instances (Ho and Szeto, 2017). Hence, most studies have focused on developing inexact methods to obtain good solutions within reasonable computational time. Heuristics such as the PILOT construction procedure (Rainer-Harbach et al., 2015), the two-phase heuristic based on linear programming by You (2019), and cluster-first, route-second (CFRS) decomposition (Schuijbroek et al., 2017), as well as metaheuristics such as iterated tabu search (Ho and Szeto, 2014), chemical reaction optimization (Szeto et al., 2016), and hybrid genetic search (Hu et al., 2025), have been proposed for the BRP. Among these studies, CFRS decomposition, which partitions the station network into clusters and solves routing subproblems within or across them, has been widely adopted for large-scale instances.

The CFRS literature can be organized by how clusters are matched to vehicles at the routing phase, which further shapes the clustering criterion and the routing-phase objective. Several studies assign each cluster to exactly one vehicle, reducing the multi-vehicle BRP to independent single-vehicle subproblems. Schuijbroek et al. (2017) derived station-level inventory bounds from a non-stationary queueing model and used these bounds to define service-level feasibility in a set-partitioning mixed-integer programming (MIP) model. Both phases minimize makespan: the clustering phase approximates it via a maximum spanning star, and the routing phase computes it exactly within each cluster. Liu et al. (2025) combined distance-minimizing clustering with inventory balance constraints and solved a single-vehicle routing MIP per cluster that minimized travel distance and a penalty on unsatisfied rebalancing quantities.

Other studies adopt decompositions in which each vehicle serves stations across multiple clusters rather than being confined to one. Forma et al. (2015) clustered stations using a savings heuristic that evaluated the reduction in UDF achieved by merging pairs of clusters. The routing phase determined both the sequence of clusters visited by each vehicle and the loading quantities at individual stations, minimizing a weighted sum of UDF and travel distance. Lv et al. (2020) clustered stations by a distance-based greedy chain heuristic with destroy-and-repair refinement, and required each cluster to be self-sufficient: the net repositioning demand of every cluster had to lie within a given threshold. An adaptive variable neighbourhood search then solved intra- and inter-cluster routing to minimize traveling cost, inventory cost, and vehicle fixed cost, under a multi-depot BSS setting. Lv et al. (2024) adopted a similar structure in a single-depot setting, using feature-correlation-reinforced clustering that combines geographic proximity and demand complementarity, and minimized CO$_2$ emissions and unmet demand in a bi-objective routing formulation.

Some studies further introduce satellite or hub stations that mediate between the depot and ordinary stations, forming a two-echelon routing structure. Lv et al. (2022) employed a fuzzy clustering method based on distance and demand correlation to designate satellite stations within each cluster. Their routing phase spanned three operational stages and minimized traveling cost, inventory cost, and a penalty on unfulfilled demand. Huang et al. (2020) instead selected hub stations via a p-hub location model that minimized the weighted sum of user walking cost from spokes to hubs and inter-hub vehicle routing cost. The routing subproblems rooted at each hub minimized a weighted sum of unmet demand and routing cost. Table 1.1 compares these CFRS-based BRP studies by their routing-phase objective, clustering criterion, and problem scale.

\begin{table}[htbp]
\centering
\caption{Comparison of CFRS approaches for the static BRP.}
\label{tab:cfrs-comparison}
\scriptsize
\setlength{\tabcolsep}{3pt}
\begin{adjustbox}{max width=\linewidth}
\begin{tabular}{p{2.3cm}p{2.35cm}p{3.2cm}p{3.7cm}ccc}
\toprule
References & Decomposition structure & Routing-phase objective & Clustering criterion & UDF-aware clustering (b) & Time budget in clustering (c) & Max \# stations (d) \\
\midrule
Forma et al. (2015) & Multi-cluster touring & Minimizing weighted sum of UDF and travel distance & Maximizing UDF reduction by cluster merging & \checkmark &  & 200 \\
Schuijbroek et al. (2017) & One cluster per vehicle & Minimizing makespan & Minimizing makespan with service-level feasibility under routing cost approximation (a) &  & \checkmark & 135 \\
Huang et al. (2020) & Hub-and-spoke & Minimizing weighted sum of unmet demand and routing cost & Minimizing demand-weighted walking and inter-hub routing cost &  &  & 200 \\
Lv et al. (2020) & Multi-cluster touring & Minimizing total travel, inventory, and vehicle fixed costs & Self-sufficiency of net repositioning demand in each cluster &  &  & 519 \\
Lv et al. (2022) & Two-echelon (satellite) & Minimizing travel, inventory, and unfulfilled-demand penalty & Distance and demand correlation &  &  & 519 \\
Lv et al. (2024) & Multi-cluster touring & Minimizing CO$_2$ emissions and unmet demand & Distance and demand correlation &  &  & 519 \\
Liu et al. (2025) & One cluster per vehicle & Minimizing travel distance with rebalancing penalty & Distance with inventory balance constraints &  &  & 615 \\
This paper & One cluster per vehicle & Maximizing UDF reduction & Maximizing estimated UDF reduction under time budget (a) & \checkmark & \checkmark & 1,000 \\
\bottomrule
\end{tabular}
\end{adjustbox}
\begin{tablenotes}\footnotesize
\item (a) The within-cluster routing cost is approximated by the maximum spanning star (Schuijbroek et al., 2017).
\item (b) Whether the clustering criterion is UDF-aware, i.e., directly evaluates UDF reduction.
\item (c) Whether the clustering phase accounts for vehicle routing time feasibility.
\item (d) Number of stations in the largest tested instance.
\end{tablenotes}
\end{table}

Two complementary gaps can be identified from Table 1.1. First, CFRS-based studies predominantly adopt distance- or demand-based clustering criteria that are not designed to capture UDF reduction. Even Liu et al. (2025), who exploit UDF convexity to derive station-level rebalancing targets, eventually adopt deviation-based penalties in both the clustering and routing phases. As UDF is a nonlinear convex function of station inventory (Raviv and Kolka, 2013), clusters that appear well-formed under such metrics may not correspond to clusters that yield substantial UDF reduction. Second, the non-CFRS studies that adopt UDF (Ho and Szeto, 2017; Raviv et al., 2013; Zhang et al., 2017) do not employ a clustering-based decomposition. Taken together, UDF-based optimization and CFRS decomposition have so far been studied largely in separate lines of the BRP literature.

The sole study at the intersection of these two lines is Forma et al. (2015), who evaluated UDF reduction within a savings heuristic for clustering. Their saving score, however, depended only on the inventory-based cost reduction of merging two clusters, with geographic proximity enforced separately through a pairwise diameter constraint and vehicle routing time feasibility deferred to the subsequent routing phases. The largest instances they tested contained 200 stations. To the best of our knowledge, no existing CFRS study formulates the clustering phase as an optimization model that jointly captures UDF reduction and routing time feasibility.

This paper proposes a CFRS framework for the static BRP that retains UDF as the objective in both the clustering and routing phases. The clustering phase is formulated as a mixed-integer linear program (MILP) that maximizes the estimated total UDF reduction across clusters, subject to vehicle time budget constraints. Within-cluster routing cost is approximated by the maximum spanning star, following Schuijbroek et al. (2017). The model jointly determines station-to-cluster assignments, approximate within-cluster bike transfers, and the resulting post-transfer inventories, embedding the nonlinear relationship between inventory and expected service failures into the clustering objective. Because this MILP becomes intractable for large networks, we develop the Hybrid Genetic Search for UDF-aware Clustering (HGS-UC) algorithm to solve it at city scale.

The main contributions of this paper are as follows:

\begin{enumerate}
\item This paper proposes a CFRS framework for the static BRP in which the clustering phase is formulated as a MILP that jointly optimizes UDF reduction and routing time feasibility. To the best of our knowledge, no prior CFRS framework has done so.
\item This paper develops the HGS-UC algorithm, which combines population-based genetic search with an Activated Transfer Selection (ATS) heuristic for efficient cluster evaluation.
\item Computational experiments on New York Citi Bike instances show that the proposed framework achieves solution quality comparable to an exact BRP solver on instances of 50 and 100 stations and achieves higher UDF reduction than five baseline clustering methods on instances of up to 1,000 stations.
\end{enumerate}

The remainder of this paper is organized as follows. Section 2 defines the multi-vehicle static BRP and presents the mathematical formulation. Section 3 introduces the CFRS framework and formulates the UDF-aware clustering model. Section 4 presents the HGS-UC algorithm. Section 5 reports computational experiments and results. Section 6 offers concluding remarks.

\section{Model formulation}

This section defines the multi-vehicle static BRP (MV-BRP) studied in this paper and presents its mixed-integer linear programming formulation.

We consider a BSS consisting of a single depot and a set of stations. Each station is characterized by its capacity, its initial bike inventory, and a user dissatisfaction function. The user dissatisfaction function represents the expected shortage of bikes and docks incurred by station users as a function of the station’s inventory after repositioning is carried out (Raviv and Kolka, 2013). Travel times between all pairs of nodes are known. The initial inventory at each station is assumed to be known, as modern BSSs provide the operator with real-time inventory data. The depot is assumed to hold a sufficient inventory of bikes and to have sufficient capacity to store returned bikes.

A fleet of homogeneous vehicles with limited capacity is used to redistribute bikes across the stations. A station whose initial inventory exceeds the level that minimizes its UDF is called a surplus station, and a station whose initial inventory falls below that level is called a deficit station. Vehicles depart from the depot simultaneously, collect bikes from surplus stations and deliver bikes to deficit stations, and finally return to the depot within a given time budget. Loading and unloading each bike takes a fixed amount of time. Repositioning is performed overnight, when user activity is negligible. Station inventories change only through vehicle operations. We assume that each station may be visited by at most one vehicle. The inventory at each station after repositioning sets the starting condition for the following operating day. The objective is to determine the route of each vehicle and the number of bikes loaded and unloaded at each visited station to maximize the total reduction in user dissatisfaction across all stations.

We adopt the user dissatisfaction function introduced by Raviv and Kolka (2013). For each station i, ${{f}_{i}}\left( s \right)$ measures the expected number of failed rentals and returns over the planning horizon as a function of the station’s inventory $s$ after the repositioning, and is convex over the integer inventory levels $0,1,\ldots ,{{c}_{i}}$. Convexity permits an exact piecewise-linear representation that can be embedded in a linear program. The UDF values at all inventory levels are precomputed and treated as input to the optimization model. We refer the reader to Raviv and Kolka (2013) for the computational procedure.

\subsection{Notation}

\begin{longtable}{@{}p{0.28\linewidth}p{0.66\linewidth}@{}}
\caption{Notation for the MV-BRP model.}\label{tab:mvbrp-notation}\\
\toprule
Symbol & Description \\
\midrule
\endfirsthead
\toprule
Symbol & Description \\
\midrule
\endhead
\bottomrule
\endfoot
\multicolumn{2}{@{}l}{\textit{Sets}} \\
$V$ & Set of stations \\
${{V}_{0}}$ & Set of all nodes, including the stations and the depot \\
$K$ & Set of vehicles \\
\multicolumn{2}{@{}l}{\textit{Parameters}} \\
${{c}_{i}}$ & Capacity of station $i$ \\
$s_{i}^{0}$ & Initial inventory at station $i$ \\
$Q$ & Vehicle capacity \\
$T$ & Time budget per vehicle \\
${{t}_{ij}}$ & Travel time from node $i$ to node $j$ \\
${{t}^{\text{load}}}$ & Time to load or unload one bike \\
${{a}_{i\ell }}$ & Intercept coefficient in the linearization of ${{f}_{i}}$ at inventory level $\ell$ \\
${{b}_{i\ell }}$ & Slope coefficient in the linearization of ${{f}_{i}}$ at inventory level $\ell$ \\
\multicolumn{2}{@{}l}{\textit{Decision variables}} \\
${{x}_{ijk}}$ & Binary variable that equals one if vehicle $k$ travels directly from node $i$ to node $j$, and zero otherwise \\
${{y}_{ijk}}$ & Number of bikes on vehicle $k$ when it travels directly from node $i$ to node $j$ \\
$y_{ik}^{\text{load}}$ & Number of bikes loaded onto vehicle $k$ at node $i$ \\
$y_{ik}^{\text{unload}}$ & Number of bikes unloaded from vehicle $k$ at node $i$ \\
${{s}_{i}}$ & Inventory at station $i$ after the repositioning \\
${{g}_{i}}$ & Linearized user dissatisfaction at station $i$ after the repositioning \\
\multicolumn{2}{@{}l}{\textit{Function}} \\
${{f}_{i}}(s)$ & User dissatisfaction function at station $i$ \\
\end{longtable}

\subsection{Arc-indexed formulation}

For the sake of completeness, we present the multi-vehicle static BRP as a mixed-integer linear program, building on the arc-indexed formulation of Raviv et al. (2013). We specialize to a homogeneous fleet with shared capacity $Q$ and equal loading and unloading times $t^{\text{load}}$, and restrict each station to be visited by at most one vehicle. The formulation is

\[
\max \sum_{i \in V} \bigl( f_i(s_i^0) - g_i \bigr)
\]

subject to

\[
s_i = s_i^0 - \sum_{k \in K} (y_{ik}^{\text{load}} - y_{ik}^{\text{unload}}) \quad \forall i \in V
\]

\[
y_{ik}^{\text{load}}-y_{ik}^{\text{unload}}=\sum\limits_{\begin{smallmatrix}
j\in {{V}_{0}} \\
j\ne i
\end{smallmatrix}}{{{y}_{ijk}}}-\sum\limits_{\begin{smallmatrix}
j\in {{V}_{0}} \\
j\ne i
\end{smallmatrix}}{{{y}_{jik}}}\quad \forall i\in {{V}_{0}},\ k\in K
\]

\[
{{y}_{ijk}}\le Q{{x}_{ijk}}\quad \forall i,j\in {{V}_{0}}(i\ne j),\ k\in K
\]

\[
\sum\limits_{\begin{smallmatrix}
j\in {{V}_{0}} \\
j\ne i
\end{smallmatrix}}{{{x}_{ijk}}}=\sum\limits_{\begin{smallmatrix}
j\in {{V}_{0}} \\
j\ne i
\end{smallmatrix}}{{{x}_{jik}}}\quad \forall i\in {{V}_{0}},\ k\in K
\]

\[
\sum_{k \in K} \sum_{\substack{j \in V_0 \\ j \neq i}} x_{jik} \leq 1 \quad \forall i \in V
\]

\[
\sum_{k \in K} y_{ik}^{\text{load}} \leq s_i^0 \quad \forall i \in V
\]

\[
\sum\limits_{k\in K}{y_{ik}^{\text{unload}}}\le {{c}_{i}}-s_{i}^{0}\quad \forall i\in V
\]

\[
\sum_{i \in V_0} (y_{ik}^{\text{load}} - y_{ik}^{\text{unload}}) = 0 \quad \forall k \in K
\]

\[
{{t}^{\text{load}}}\cdot \sum\limits_{i\in {{V}_{0}}}{(y_{ik}^{\text{load}}+y_{ik}^{\text{unload}})}+\sum\limits_{\begin{smallmatrix}
i,j\in {{V}_{0}} \\
i\ne j
\end{smallmatrix}}{{{t}_{ij}}}{{x}_{ijk}}\le T\quad \forall k\in K
\]

\[
\sum\limits_{j\in V}{{{x}_{0jk}}}=1\quad \forall k\in K
\]

\[
\sum\limits_{j\in V}{{{x}_{j0k}}}=1\quad \forall k\in K
\]

\[
y_{ik}^{\text{load}}\le s_{i}^{0}{{\sum }_{\begin{matrix}
j\in {{V}_{0}}  \\
j\ne i  \\
\end{matrix}}}{{x}_{jik}}\quad \forall i\in V,\ k\in K
\]

\[
y_{ik}^{\text{unload}}\le ({{c}_{i}}-s_{i}^{0}){{\sum }_{\begin{matrix}
j\in {{V}_{0}}  \\
j\ne i  \\
\end{matrix}}}{{x}_{jik}}\quad \forall i\in V,\ k\in K
\]

\[
{{g}_{i}}\ge {{a}_{i\ell }}+{{b}_{i\ell }}{{s}_{i}}\quad \forall i\in V,\ \ell =0,1,\ldots ,{{c}_{i}}-1
\]

\[
\sum\limits_{i\in S}{\sum\limits_{\begin{smallmatrix}
j\in S \\
j\ne i
\end{smallmatrix}}{{{x}_{ijk}}}}\le |S|-1\quad \forall S\subseteq V,\ |S|\ge 2,\ \forall k\in K
\]

\[
{{x}_{ijk}}\in \{0,1\}\quad \forall i,j\in {{V}_{0}}(i\ne j),k\in K
\]

\[
y_{ik}^{\text{load}}\ge 0,\ \text{integer}\quad \forall i\in {{V}_{0}},\ k\in K
\]

\[
y_{ik}^{\text{load}}\ge 0,\ \text{integer}\quad \forall i\in {{V}_{0}},\ k\in K
\]

\[
{{y}_{ijk}}\ge 0\quad \forall i,j\in {{V}_{0}}(i\ne j),k\in K
\]

\[
{{s}_{i}}\ge 0\quad \forall i\in V
\]

\[
{{s}_{i}}\le {{c}_{i}}\quad \forall i\in V
\]

\[
{{g}_{i}}\ge 0\quad \forall i\in V
\]

The objective function  maximizes the total reduction in user dissatisfaction across all stations: for each station $i$, the term $f_i(s_i^0) - g_i$ measures the improvement from the initial UDF level to the level after the repositioning. Constraints  define the post-repositioning inventory at each station as the initial inventory adjusted by all loading and unloading operations. Constraints  ensure conservation of bikes on each vehicle: the net loading at a node equals the change in vehicle load. Constraints  limit the load on each vehicle to its capacity and set the load to zero on arcs not traversed. Constraints  are the vehicle flow conservation equations. Constraints  stipulate that each station can be visited by at most one vehicle across the fleet.

Constraints  and  bound the total loading and unloading at each station by the available inventory and remaining capacity, respectively. Constraints  require that all bikes collected by each vehicle are eventually delivered. Constraints  limit the total operating time of each vehicle, including travel, loading, and unloading, to the time budget. Constraints  and  require each vehicle to depart from and return to the depot exactly once. Constraints  and  tighten the loading and unloading bounds per vehicle by conditioning on whether that vehicle visits the station. These valid inequalities strengthen the LP relaxation (Raviv et al., 2013).

The UDF at each station is linearized following Raviv et al. (2013). Constraints  are the linear constraints that support the convex function ${{f}_{i}}$. For each station $i$ and inventory level $\ell =0,1,\ldots ,{{c}_{i}}-1$, the slope and intercept are ${{b}_{i\ell }}={{f}_{i}}\left( \ell +1 \right)-{{f}_{i}}\left( \ell  \right)$ and ${{a}_{i\ell }}={{f}_{i}}\left( \ell  \right)-{{b}_{i\ell }}\cdot \ell $, respectively. Because $f_i$ is convex and defined over integer values, this linearization is exact (Raviv et al., 2013). Constraints  are the Dantzig-Fulkerson-Johnson (DFJ) subtour elimination constraints (Dantzig et al., 1954), which prevent any subset of stations from forming a disconnected tour. Constraints - define the variable domains.

\section{The CFRS framework}

This section presents the CFRS framework for the multi-vehicle BRP. The framework decomposes the problem into a clustering phase, which partitions the station set $V$ into $\left|K\right|$ disjoint groups (one per vehicle) by maximizing the total estimated UDF reduction subject to a per-cluster time budget, and a routing phase, which solves a single-vehicle instance of the BRP within each cluster.

\begin{figure}[htbp]
\centering
\fbox{\parbox{0.86\linewidth}{\centering Figure 3.1 placeholder. Original image was not present in the txt export.}}
\caption{Illustration of the CFRS framework on a BSS network with one depot, nine stations, and $\left| K \right|=3$ vehicles (clusters).}
\end{figure}

Fig. 3.1 illustrates the framework on a small example with one depot and nine stations partitioned into $\left| K \right|=3$ clusters. In the figure, square node 0 denotes the depot. Red and blue circles denote surplus stations ($s_{i}^{0}>s_{i}^{*}$) and deficit stations ($s_{i}^{0}<s_{i}^{*}$), respectively. Bars next to each station show the current inventory $s_{i}^{0}$ relative to the optimal level $s_{i}^{*}$. Phase 1 partitions the stations into clusters and determines approximate within-cluster transfers between surplus and deficit stations. Phase 2 solves a single-vehicle BRP within each cluster. The system-wide UDF reduction is the sum of the per-cluster reductions.

Section 3.1 formulates the clustering phase as a mixed-integer linear program. Section 3.2 describes the routing phase.

\subsection{UDF-aware clustering model}

We formulate the clustering phase as a mixed-integer linear program that jointly assigns stations to clusters and determines approximate transfers between surplus and deficit stations within each cluster, subject to a per-cluster time budget and cluster size limits. The notation from Table 2.1 carries over. Table 3.1 lists the additional sets, parameters, and decision variables.

\subsubsection{Notation}

\begin{longtable}{@{}p{0.28\linewidth}p{0.66\linewidth}@{}}
\caption{Additional notation for the clustering model.}\label{tab:clustering-notation}\\
\toprule
Symbol & Description \\
\midrule
\endfirsthead
\toprule
Symbol & Description \\
\midrule
\endhead
\bottomrule
\endfoot
\multicolumn{2}{@{}l}{\textit{Sets}} \\
${{V}^{+}}$ & Set of surplus stations, ${{V}^{+}}=\{i\in V \mid s_{i}^{0}>s_{i}^{*}\}$ \\
${{V}^{-}}$ & Set of deficit stations, ${{V}^{-}}=\{i\in V \mid s_{i}^{0}<s_{i}^{*}\}$ \\
\multicolumn{2}{@{}l}{\textit{Parameters}} \\
$s_{i}^{*}$ & Inventory level that minimizes ${{f}_{i}}$ at station $i$ \\
${{U}_{ij}}$ & Maximum number of bikes transferable from station $i$ to station $j$, ${{U}_{ij}}=\min \left\{ s_{i}^{0}-s_{i}^{*},s_{j}^{*}-s_{j}^{0} \right\}$ if $i\in {{V}^{+}}$ and $j\in {{V}^{-}}$, 0 otherwise \\
${{N}_{\min }}$ & Minimum number of stations per cluster \\
${{N}_{\max }}$ & Maximum number of stations per cluster \\
\multicolumn{2}{@{}l}{\textit{Decision variables}} \\
${{z}_{ik}}$ & Binary variable that equals one if station $i$ is assigned to cluster $k$, and zero otherwise \\
${{w}_{ijk}}$ & Number of bikes transferred from station $i$ to station $j$ within cluster $k$ \\
${{\delta }_{ik}}$ & Binary variable that equals one if station $i$ participates in at least one transfer within cluster $k$, and zero otherwise \\
$\tau _{k}^{\text{star}}$ & Maximum spanning star component of the estimated travel time for cluster $k$ \\
$\tau _{k}^{\text{depot}}$ & Depot round-trip component of the estimated travel time for cluster $k$ \\
$\tilde{s}_i$ & Estimated inventory at station $i$ after the bike transfers \\
$\tilde{g}_i$ & Estimated linearized UDF at station $i$ after the bike transfers \\
\end{longtable}

\subsubsection{MILP formulation}

The formulation is

\[
\max \sum\limits_{i\in V}{\left( {{f}_{i}}(s_{i}^{0})-{{{\tilde{g}}}_{i}} \right)}
\]

subject to

\[
\sum\limits_{k\in K}{{{z}_{ik}}}=1\quad \forall i\in V
\]

\[
{{w}_{ijk}}\le {{U}_{ij}}{{z}_{ik}}\quad \forall i\in {{V}^{+}},j\in {{V}^{-}},k\in K
\]

\[
{{w}_{ijk}}\le {{U}_{ij}}{{z}_{jk}}\quad \forall i\in {{V}^{+}},j\in {{V}^{-}},k\in K
\]

\[
\sum\limits_{j\in {{V}^{-}}}{\sum\limits_{k\in K}{{{w}_{ijk}}}}\le s_{i}^{0}-s_{i}^{*}\quad \forall i\in {{V}^{+}}
\]

\[
\sum\limits_{i\in {{V}^{+}}}{\sum\limits_{k\in K}{{{w}_{ijk}}}}\le s_{j}^{*}-s_{j}^{0}\quad \forall j\in {{V}^{-}}
\]

\[
{{\delta }_{ik}}\le {{z}_{ik}}\quad \forall i\in V,k\in K
\]

\[
\sum\limits_{j\in {{V}^{-}}}{{{w}_{ijk}}}\le \left( s_{i}^{0}-s_{i}^{*} \right){{\delta }_{ik}}\quad \forall i\in {{V}^{+}},k\in K
\]

\[
\sum\limits_{j\in {{V}^{-}}}{{{w}_{ijk}}}\ge {{\delta }_{ik}}\quad \forall i\in {{V}^{+}},k\in K
\]

\[
\sum\limits_{i\in {{V}^{+}}}{{{w}_{ijk}}}\le \left( s_{j}^{*}-s_{j}^{0} \right){{\delta }_{jk}}\quad \forall j\in {{V}^{-}},k\in K
\]

\[
\sum\limits_{i\in {{V}^{+}}}{{{w}_{ijk}}}\ge {{\delta }_{jk}}\quad \forall j\in {{V}^{-}},k\in K
\]

\[
\tau _{k}^{\text{star}}\ge \sum\limits_{i\in V}{\max \left\{ {{t}_{ij}},{{t}_{ji}} \right\}}\cdot \left( {{\delta }_{ik}}+{{\delta }_{jk}}-1 \right)\quad \forall j\in V,k\in K
\]

\[
\tau _{k}^{\text{depot}}\ge 2\max \left\{ {{t}_{0j}},{{t}_{j0}} \right\}\cdot {{\delta }_{jk}}\quad \forall j\in V,k\in K
\]

\[
\tau _{k}^{\text{star}}+\tau _{k}^{\text{depot}}+2{{t}^{\text{load}}}\sum\limits_{i\in {{V}^{+}}}{\sum\limits_{j\in {{V}^{-}}}{{{w}_{ijk}}}}\le T\quad \forall k\in K
\]

\[
\sum\limits_{i\in V}{{{z}_{ik}}}\ge {{N}_{\min }}\quad \forall k\in K
\]

\[
\sum\limits_{i\in V}{{{z}_{ik}}}\le {{N}_{\max }}\quad \forall k\in K
\]

\[
\tilde{s}_i=s_{i}^{0}-\sum\limits_{j\in {{V}^{-}}}{\sum\limits_{k\in K}{{{w}_{ijk}}}}\quad \forall i\in {{V}^{+}}
\]

\[
\tilde{s}_j=s_{j}^{0}+\sum\limits_{i\in {{V}^{+}}}{\sum\limits_{k\in K}{{{w}_{ijk}}}}\quad \forall j\in {{V}^{-}}
\]

\[
{{\tilde{s}}_{i}}=s_{i}^{0}\quad \forall i\in V\backslash \left( {{V}^{+}}\cup {{V}^{-}} \right)
\]

\[
{{\tilde{g}}_{i}}\ge {{a}_{i\ell }}+{{b}_{i\ell }}{{\tilde{s}}_{i}}\quad \forall i\in V,\ell =0,1,\ldots ,{{c}_{i}}-1
\]

\[
{{z}_{ik}}\in \{0,1\}\quad \forall i\in V,k\in K
\]

\[
{{\delta }_{ik}}\in \{0,1\}\quad \forall i\in V,k\in K
\]

\[
{{w}_{ijk}}\ge 0,\ \text{integer}\quad \forall i\in {{V}^{+}},j\in {{V}^{-}},k\in K
\]

\[
\tau _{k}^{\text{star}}\ge 0\quad \forall k\in K
\]

\[
\tau _{k}^{\text{depot}}\ge 0\quad \forall k\in K
\]

\[
\tilde{s}_i \ge 0\quad \forall i\in V
\]

\[
{{\tilde{s}}_{i}}\le {{c}_{i}}\quad \forall i\in V
\]

\[
\tilde{g}_i \ge 0\quad \forall i\in V
\]

The objective function  has the same form as  in Section 2.2, with $\tilde{g}_i$ replacing $g_i$ to indicate that the UDF is evaluated at the estimated inventory after the transfers rather than at the actual repositioning outcome. Constraints  ensure that each station is assigned to exactly one cluster. Constraints  and  ensure that a transfer between two stations can take place only if both are assigned to the same cluster.

Constraints  and  bound the total transfers out of each surplus station and into each deficit station by the available surplus and deficit, respectively. Constraints – tie the activation variables to the assignment and transfer variables, so that ${{\delta }_{ik}}=1$ if and only if station $i$ is assigned to cluster $k$ and participates in at least one transfer within it. Balanced stations ($s_{i}^{0}=s_{i}^{*}$) are never activated, as they have neither surplus to offer nor deficit to fill.

Constraints – approximate the within-cluster routing time as the sum of a maximum spanning star (MSS) component and a depot round-trip component, following Schuijbroek et al. (2017). Specifically, it is defined as the maximum over all candidate centers of the sum of distances from the center to all other active stations in the cluster, which is an upper bound on the shortest Hamiltonian path through the cluster and admits a compact linear representation. Constraints  encode this maximum through the weight $\left( {{\delta }_{ik}}+{{\delta }_{jk}}-1 \right)$: when $j$ is an active center (${{\delta }_{jk}}=1$), it reduces to ${{\delta }_{ik}}$ and the right-hand side becomes the sum of distances from j to all other active stations in cluster k. Otherwise, every term is non-positive and the inequality is deactivated. Constraints  account for the depot round-trip through the farthest active station. Constraints  enforce the time budget, where the factor 2 on ${{t}^{\text{load}}}$ reflects the loading time at the source and the unloading time at the destination of each transferred bike.

When the travel time matrix is asymmetric, constraints (35) and (36) use the symmetrized value $\text{max}\left\{ {{t}_{ij}},{{t}_{ji}} \right\}$ on each station pair rather than a single directed value. Since the within-cluster route is not yet determined at the clustering stage, the orientation in which any edge would be traversed is unknown. Taking the maximum of the two directions yields a symmetric surrogate that upper-bounds the true routing cost in either orientation, so a cluster satisfying the time budget under this surrogate remains feasible once the actual route is computed.

Constraints  and  bound the cluster sizes between $N_{\min}$ and $N_{\max}$. Constraints – define the estimated inventory at each station as the initial inventory adjusted by the within-cluster transfers. Constraints  linearize the estimated UDF using the same piecewise-linear technique as constraints  in Section 2.2. The representation is exact because $f_i$ is convex over integer inventory levels. Constraints – define the variable domains.

Compared with the savings heuristic of Forma et al. (2015), which also incorporates UDF into clustering but evaluates each cluster’s UDF reduction by optimally redistributing its total bike count under the assumption of costless within-cluster transfers, our formulation differs in two structural aspects. First, transfer quantities $w_{ijk}$ between individual surplus–deficit pairs are determined explicitly, bounded by each station's available surplus or deficit, so that UDF is evaluated at station-level rather than cluster-level inventories. Second, these transfers are subject to the per-cluster time budget through the MSS approximation in –, so that geographic proximity and routing feasibility enter the clustering decision directly rather than through separate proximity constraints.

The MSS is known to consistently overestimate the actual routing cost (Schuijbroek et al., 2017), so the time budget constraints  are enforced against an upper bound rather than the exact shortest Hamiltonian path. This keeps the clustering model linear and free of exact routing variables, at the cost of excluding some clusters whose exact single-vehicle route would in fact fit within the time budget. Nonetheless, the clustering MILP remains intractable for city-scale instances, motivating the metaheuristic developed in Section 4.

\subsection{Routing phase}

Given the partition produced by the clustering phase, each cluster is solved as an independent single-vehicle instance of the BRP formulated in Section 2.2, restricted to the stations assigned to that cluster together with the depot. Only the station-to-cluster assignment is passed to the routing phase. The approximate transfers $w_{ijk}$ from the clustering phase are discarded, and each subproblem is solved from scratch on its own UDF curves and travel time submatrix.

\section{Hybrid genetic search for UDF-aware clustering}

This section presents the Hybrid Genetic Search for UDF-aware Clustering (HGS-UC), a metaheuristic for the clustering model formulated in Section 3.1. HGS-UC adapts the Hybrid Genetic Search framework of Vidal et al. (2012), refined in the open-source implementation of Vidal (2022).

\subsection{Solution representation}

A clustering solution is represented by $\phi =({{C}_{1}},{{C}_{2}},\ldots ,{{C}_{\left| K \right|}})$, where $K$ is the homogeneous vehicle set and ${{C}_{k}}$ denotes the set of stations assigned to cluster $k$. The clusters form a partition of $V$, that is $\bigcup\limits_{k=1}^{\left| K \right|}{{{C}_{k}}}=V\text{ and }{{C}_{k}}\cap {{C}_{{{k}'}}}=\varnothing ,\text{ }\forall k\ne {k}'.$ The depot is not assigned to any cluster and is not represented in the solution. Since each cluster is served by exactly one vehicle, the terms cluster and vehicle are used interchangeably throughout this section.

For example, given $\left| V \right|=9$ stations and $\left| K \right|=3$ vehicles, one feasible solution is $\phi =(\{2,5,7\},\{1,4,9\},\{3,6,8\})$, in which stations 2, 5, and 7 are assigned to vehicle 1, stations 1, 4, and 9 to vehicle 2, and stations 3, 6, and 8 to vehicle 3.

\subsection{Algorithm overview}

The algorithmic procedure of HGS-UC is given in Algorithm 1. The algorithm evolves a population $P$ of clustering solutions. The objective value of a solution $\phi$ is denoted as $F(\phi)$\footnote{Throughout this section, $F$ denotes the objective value of the optimization problem, while $f$ retains its meaning as the UDF defined in Sections 2 and 3.}.

The algorithm first generates an initial population $P$ (Line 1) and identifies the best solution $\phi^*$ (Line 2). It then records the current best objective value as ${{F}_{\text{prev}}}$ and initializes the non-improvement counters $ite{{r}_{\text{NI}}}$ and $ite{{r}_{\text{div}}}$ to zero (Line 3). The population then iteratively evolves through successive operations (Lines 4–25). In each generation, two parents are selected by binary tournament based on their biased fitness values (Line 5) and recombined with swapped parent roles to produce two offspring via crossover (Lines 6–11), each undergoing cluster size repair and education before entering the population. Survivor selection is applied if the population exceeds the maximum size (Lines 12–14). After updating $\phi$ (Line 15), the counters $ite{{r}_{\text{NI}}}$ and $ite{{r}_{\text{div}}}$ are reset to 0 if $F({{\phi }^{*}})$ improves over ${{F}_{\text{prev}}}$, and are incremented by 1 otherwise (Lines 16–20). A diversification phase is activated every $I{{t}_{\text{div}}}$ iterations without improvement (Lines 21–24), retaining a small set of solutions and regenerating the remainder through the initialization procedure. The algorithm terminates and returns the best-known solution $\phi^*$ and its objective value $F({{\phi }^{*}})$ (Line 26) when the number of consecutive non-improving iterations $ite{{r}_{\text{NI}}}$ reaches the limit $I{{t}_{\text{NI}}}$ or the elapsed time reaches ${{T}_{\max }}$.

\begin{algorithm}[htbp]
\caption{HGS-UC}
\label{alg:hgs-uc}
\begin{algorithmic}[1]
\State $P\leftarrow$ InitPopulation
\State $\phi^{*}\leftarrow \arg\max_{\phi\in P}F(\phi)$
\State Set $F_{\text{prev}}\leftarrow F(\phi^{*})$, $iter_{\text{NI}}\leftarrow 0$, and $iter_{\text{div}}\leftarrow 0$
\While{$iter_{\text{NI}}<It_{\text{NI}}$ and elapsed time $<T_{\max}$}
\State Select $\phi_{1}$ and $\phi_{2}$ from $P$ by binary tournament based on biased fitness
\For{$(\phi_A,\phi_B)\in\{(\phi_1,\phi_2),(\phi_2,\phi_1)\}$}
\State $\phi\leftarrow$ Crossover$(\phi_A,\phi_B)$
\State ClusterSizeRepair$(\phi)$
\State Educate$(\phi)$
\State Add $\phi$ to $P$
\EndFor
\If{the size of $P$ exceeds the maximum population size}
\State SelectSurvivors$(P)$
\EndIf
\State Update $\phi^{*}$ if a better solution exists in $P$
\If{$F(\phi^{*})>F_{\text{prev}}$}
\State Set $F_{\text{prev}}\leftarrow F(\phi^{*})$, $iter_{\text{NI}}\leftarrow 0$, and $iter_{\text{div}}\leftarrow 0$
\Else
\State Set $iter_{\text{NI}}\leftarrow iter_{\text{NI}}+1$ and $iter_{\text{div}}\leftarrow iter_{\text{div}}+1$
\EndIf
\If{$iter_{\text{div}}\ge It_{\text{div}}$}
\State Diversify$(P)$
\State Set $iter_{\text{div}}\leftarrow 0$
\EndIf
\EndWhile
\State \Return $\phi^{*},F(\phi^{*})$
\end{algorithmic}
\end{algorithm}

Unlike the original HGS framework, which maintains separate subpopulations of feasible and infeasible solutions, HGS-UC operates on a single feasible population. Feasibility is enforced by the cluster size repair (Line 8) invoked before each individual enters the population.

The solution evaluation procedure is detailed in Section 4.3. Section 4.4 describes solution construction. Section 4.5 describes crossover. Cluster size repair, the education operators, and population management are presented in Sections 4.6, 4.7, and 4.8, respectively.

\subsection{Solution evaluation}

To evaluate a solution $\phi =({{C}_{1}},\ldots ,{{C}_{\left| K \right|}})$, the algorithm computes the total estimated UDF reduction across all clusters. The UDF reduction at each station is determined by its inventory after the repositioning, which in turn depends on the bike transfers carried out within its cluster. Exact computation of the optimal within-cluster transfers and the corresponding route requires solving a single-vehicle BRP per cluster, which is impractical inside a population-based metaheuristic that evaluates many clusters per generation.

Therefore, we propose the Activated Transfer Selection (ATS) procedure, which estimates a time-feasible transfer set for each cluster. ATS scores each candidate transfer by its expected UDF improvement per unit handling time, selects transfers in decreasing score order, and terminates when the time budget is exhausted.

The ATS procedure operates on the following derived quantities, defined per cluster ${{C}_{k}}$ unless otherwise stated:

\begin{enumerate}[label=(\arabic*)]
\item Surplus and deficit quantities. For each station $i\in {{V}^{+}}$, the surplus quantity is $r_{i}^{+}=s_{i}^{0}-s_{i}^{*}$. For each station $j\in {{V}^{-}}$, the deficit quantity is $r_{j}^{-}=s_{j}^{*}-s_{j}^{0}$.
\item Edges. A directed pair $(i,j)$ with $i\in {{V}^{+}}$ and $j\in {{V}^{-}}$ is called an edge. The maximum transferable quantity on edge $(i,j)$ is $\min (r_{i}^{+},r_{j}^{-})$.
\item Transfers and transfer sets. A transfer is a triple $e=(i,j,q)$, where $(i,j)$ is an edge and $q\ge 1$ is the number of bikes moved from $i$ to $j$. A transfer set ${{\mathcal{T}}_{k}}$ collects the transfers selected for execution within cluster ${{C}_{k}}$.
\item Marginal UDF reductions. For a surplus station $i$, the $m$-th marginal UDF reduction is $\Delta _{i}^{m}={{f}_{i}}(s_{i}^{0}-m+1)-{{f}_{i}}(s_{i}^{0}-m)$ for $m=1,\ldots ,r_{i}^{+}$, representing the decrease in UDF when the $m$-th bike is removed. For a deficit station $j$, the $m$-th marginal UDF reduction is $\Delta _{j}^{m}={{f}_{j}}(s_{j}^{0}+m-1)-{{f}_{j}}(s_{j}^{0}+m)$ for $m=1,\ldots ,r_{j}^{-}$. Because ${{f}_{i}}$ is convex, the marginals are decreasing: $\Delta _{i}^{1}\ge \Delta _{i}^{2}\ge \cdots \ge 0$.
\item Edge scores. The score of edge $(i,j)$ is ${{\sigma }_{ij}}=(\Delta _{i}^{1}+\Delta _{j}^{1})/(2{{t}^{\text{load}}})$, where the denominator accounts for one loading and one unloading operation. The score measures the UDF reduction per unit handling time for the first bike transferred on the edge.
\item Activated stations. The set of stations participating in at least one accepted transfer within cluster ${{C}_{k}}$ is denoted $\mathcal{A}\subseteq {{C}_{k}}$. Each member of $\mathcal{A}$ must be visited by the vehicle, incurring routing time.
\item Routing time estimate. The routing time for a set of activated stations $\mathcal{A}$ is estimated as $\tau (\mathcal{A})=MSS(\mathcal{A})+2{{\max }_{i\in \mathcal{A}}}\max ({{t}_{0i}},{{t}_{i0}})$, where $MSS(\mathcal{A})={{\max }_{i\in \mathcal{A}}}\sum\limits_{j\in \mathcal{A},j\ne i}{\max }({{t}_{ij}},{{t}_{ji}})$ is the maximum spanning star and the second term accounts for the depot round trip. This is the same estimate used in the clustering model.
\item MSS time feasibility. A transfer set ${{\mathcal{T}}_{k}}$ is MSS time-feasible if $\tau (\mathcal{A})+2{{t}^{\text{load}}}\sum\limits_{(i,j,q)\in {{\mathcal{T}}_{k}}}{q}\le T$, where the handling term reflects that each transferred bike is loaded at its source and unloaded at its destination. Because the MSS provides an upper bound on the optimal tour length through $\mathcal{A}$, any transfer set that satisfies this condition remains feasible when the actual route is computed.
\end{enumerate}

We now present the ATS procedure in detail. The procedure takes a cluster ${{C}_{k}}$ as input and returns an MSS time-feasible transfer set ${{\mathcal{T}}_{k}}$. Each cluster is processed independently.

Step 0 (Initialization). Generate all edges within ${{C}_{k}}$ and sort them by ${{\sigma }_{ij}}$ in decreasing order. Set $\mathcal{A}\leftarrow \varnothing $ and ${{\mathcal{T}}_{k}}\leftarrow \varnothing $. Initialize the residual surplus $\tilde{r}_{i}^{+}\leftarrow r_{i}^{+}$ for all $i\in {{V}^{+}}\cap {{C}_{k}}$ and the residual deficit $\tilde{r}_{j}^{-}\leftarrow r_{j}^{-}$ for all $j\in {{V}^{-}}\cap {{C}_{k}}$.

Step 1 (First-marginal selection). Process the sorted edges sequentially. For each edge $(i,j)$, skip if $i\in \mathcal{A}$ or $j\in \mathcal{A}$. Otherwise, tentatively add $i$ and $j$ to $\mathcal{A}$, recompute $\tau (\mathcal{A})$, and compute the feasible quantity

\[
q=\min \left\{ \tilde{r}_{i}^{+},\ \tilde{r}_{j}^{-},\ \left\lfloor \frac{T-\tau (\mathcal{A})-2{{t}^{\text{load}}}\sum\limits_{({i}',{j}',{q}')\in {{\mathcal{T}}_{k}}}{{{q}'}}}{2{{t}^{\text{load}}}} \right\rfloor  \right\},
\]

If $q\ge 1$, accept transfer $(i,j,q)$ into ${{\mathcal{T}}_{k}}$ and update $\tilde{r}_{i}^{+}\leftarrow \tilde{r}_{i}^{+}-q$ and $\tilde{r}_{j}^{-}\leftarrow \tilde{r}_{j}^{-}-q$. If $q<1$, remove $i$ and $j$ from $\mathcal{A}$ and proceed to the next edge. Skipping edges whose stations are already active enforces mutual exclusivity and keeps the precomputed first-marginal scores valid throughout this step.

Step 2 (Residual capacity matching). This step fills in additional transfers that reuse stations already activated in Step 1, exploiting any remaining time budget. A station may appear in multiple transfers during this step. Each additional bike on a station is scored by its own marginal UDF reduction, so that the diminishing returns on a repeatedly used station are reflected.

At each iteration, Step 2 selects the surplus-deficit pair $(i,j)$ within ${{C}_{k}}$ with the highest score

\[
{{\tilde{\sigma }}_{ij}}=\frac{\Delta _{i}^{{{m}_{i}}+1}+\Delta _{j}^{{{m}_{j}}+1}}{2\,{{t}^{\text{load}}}},
\]

where ${{m}_{i}}$ and ${{m}_{j}}$ are the numbers of bikes already transferred from $i$ and into $j$ by previously accepted transfers in this cluster. Only pairs with ${{\tilde{\sigma }}_{ij}}>0$ are considered.

For the selected pair, any inactive station is tentatively activated. $\tau \left( \mathcal{A} \right)$ is recomputed, and the feasible quantity $q$ is determined as in Step 1. If $q\ge 1$, the transfer is accepted and all residual quantities are updated. Step 2 terminates when no surplus-deficit pair with a positive score remains, or when the remaining time budget can no longer accommodate any transfer.

After applying ATS to every cluster, the objective value of solution $\phi$ is
\[
F(\phi)=\sum_{i\in V}\left[f_i(s_i^0)-f_i(s_i^0+\nu_i)\right],
\]
where
\[
\nu_i=\begin{cases}
-\sum_{j:(i,j,q)\in\mathcal{T}} q, & \text{if } i\in V^{+},\\
\sum_{j:(j,i,q)\in\mathcal{T}} q, & \text{if } i\in V^{-},\\
0, & \text{otherwise},
\end{cases}
\]
is the net inventory change at station $i$ and $\mathcal{T}=\bigcup_{k=1}^{|K|}\mathcal{T}_k$ is the set of all accepted transfers. Fig. 4.1 illustrates the procedure on a four-station toy cluster.

\begin{figure}[htbp]
\centering
\fbox{\parbox{0.86\linewidth}{\centering Figure 4.1 placeholder. Original image was not present in the txt export.}}
\caption{Illustration of the ATS procedure on a toy cluster with surplus stations 1 and 2 and deficit stations 3 and 4. (a) Candidate edges are ranked by the first-marginal score $\sigma_{ij}$. (b) Step 1 scans the edges in decreasing $\sigma_{ij}$ order and accepts transfers (1, 3, 2) and (2, 4, 2), activating all four stations and consuming 70\% of the time budget T. (c) Step 2 revisits the previously skipped pair (1, 4) and appends a transfer of q = 1, raising time budget usage to 80\%. The accepted-transfer set $\mathcal{T}_k$ shown below each panel updates accordingly.}
\end{figure}

\subsection{Solution construction}

The algorithm constructs new solutions when building the initial population (Algorithm 1, Line 1) and when regenerating individuals during diversification (Line 22). Each solution $\phi =({{C}_{1}},\ldots ,{{C}_{\left| K \right|}})$ is constructed from scratch by applying K-medoids clustering on a composite dissimilarity measure, which combines two normalized components: geographic proximity, measured by the travel time ${{t}_{ij}}$, and UDF reduction potential, measured by the bilateral UDF reduction ${{B}_{ij}}$.

Specifically, the bilateral UDF reduction of stations i and j is defined as

\[
{{B}_{ij}}=\left\{ \begin{array}{*{35}{l}}
\max ({{f}_{i}}(s_{i}^{0})+{{f}_{j}}(s_{j}^{0})-{{f}_{i}}(s_{i}^{0}-{{q}_{ij}})-{{f}_{j}}(s_{j}^{0}+{{q}_{ij}}),\ 0) & \text{if }i\in {{V}^{+}},j\in {{V}^{-}}  \\
0 & \text{otherwise}  \\
\end{array} \right.,
\]

where ${{q}_{ij}}=\min (r_{i}^{+},r_{j}^{-})$. The dissimilarity of stations i and j is then defined as

\[
{{D}_{ij}}=\alpha {{\hat{t}}_{ij}}+(1-\alpha )(1-{{\hat{B}}_{ij}}),
\]

where ${{\hat{t}}_{ij}}$ and ${{\hat{B}}_{ij}}$ denote min–max normalizations over all off-diagonal station pairs, and $\alpha \in [0,1]$ controls the relative weight. A low ${{D}_{ij}}$ indicates that $i$ and $j$ are geographically close and offer substantial mutual UDF reduction, making them desirable cluster partners.

The construction procedure applies K-medoids clustering with K-means++ seeding (Arthur and Vassilvitskii, 2007), adapted to the K-medoids setting by sampling proportionally to the dissimilarity rather than its square. The procedure is as follows:

Step 0 (Seeding). Select the first medoid ${{m}_{1}}$ uniformly at random from $V$. For $k=2,\ldots ,\left| K \right|$, let ${{M}_{k-1}}=\{{{m}_{1}},\ldots ,{{m}_{k-1}}\}$ and ${{d}_{i}}={{\min }_{m\in {{M}_{k-1}}}}{{D}_{i,m}}$ for each $i\in V\backslash {{M}_{k-1}}$. Sample ${{m}_{k}}=i$ from $V\backslash {{M}_{k-1}}$ with probability $\frac{{{d}_{i}}}{\sum\limits_{j\in V\backslash {{M}_{k-1}}}{{{d}_{j}}}}.$

Step 1 (Assignment). Assign each station $i\in V$ to the cluster ${{C}_{k}}$ whose medoid minimizes ${{D}_{i,{{m}_{k}}}}$.

Step 2 (Medoid update). Recompute each medoid as ${{m}_{k}}=\arg {{\min }_{i\in {{C}_{k}}}}\sum\limits_{j\in {{C}_{k}}}{{{D}_{ij}}}$.

Step 3 (Convergence check). If the set of medoids has changed, return to Step 1. Otherwise, terminate and return $\phi =({{C}_{1}},\ldots ,{{C}_{\left| K \right|}})$.

Diversity across constructed solutions is enabled through the stochastic seeding in Step 0. Each invocation samples different initial medoids and produces a distinct partition.

\subsection{Crossover}

The crossover operator recombines two parent solutions to produce an offspring that inherits intra-cluster structure from one parent and inter-cluster boundaries from the other. Two offspring are produced per generation by applying the operator with swapped parent roles (Algorithm 1, Line 6). The procedure of crossover is given in Algorithm 2.

\begin{algorithm}[htbp]
\caption{Crossover$(\phi_A,\phi_B)$}
\label{alg:crossover}
\begin{algorithmic}[1]
\State $\mathcal{I}\leftarrow$ a set of $\lfloor |K|/2\rfloor$ indices sampled uniformly from $K$
\State $\mathcal{AS}\leftarrow\varnothing$ \Comment{record stations assigned to the offspring so far}
\For{$k\in\mathcal{I}$}
\State $C_k\leftarrow C_k(\phi_A)$, $m_k\leftarrow m_k(\phi_A)$, $\mathcal{AS}\leftarrow\mathcal{AS}\cup C_k$
\EndFor
\For{each $i\in V\setminus\mathcal{AS}$}
\State $k'\leftarrow$ cluster of $i$ in $\phi_B$
\If{$k'\notin\mathcal{I}$}
\State Assign $i$ to cluster $k'$
\Else
\State Assign $i$ to cluster $\arg\min_{k\in\mathcal{I}}D_{i,m_k}$
\EndIf
\EndFor
\State Reseed any empty cluster from the largest cluster
\State Recompute all medoids
\State \Return $(C_1,\ldots,C_{|K|})$
\end{algorithmic}
\end{algorithm}

Given parents ${{\phi }_{A}}$ and ${{\phi }_{B}}$, the operator first samples $\left\lfloor \left| K \right|/2 \right\rfloor $ cluster indices and copies the corresponding clusters from ${{\phi }_{A}}$ into the offspring, preserving member stations and medoids (Lines 1–5). Each remaining station $i\in V\backslash \mathcal{A}\mathcal{S}$ is then allocated using ${{\phi }_{B}}$: if ${{\phi }_{B}}$ places $i$ in a cluster ${{k}^{\prime }}$ not among the inherited indices, $i$ is assigned to ${{k}^{\prime }}$. Otherwise, $i$ is assigned to the inherited cluster whose medoid has the smallest dissimilarity ${{D}_{i,{{m}_{k}}}}$ (Lines 6–13). Any cluster that remains empty after reassignment is reseeded from the largest cluster, and all medoids are recomputed (Lines 14–15). Fig. 4.2 walks through a nine-station example with $\left|K\right|=4$ that triggers both the direct copy and the fallback assignment.

Inheriting complete clusters preserves the intra-cluster groupings refined by the donating parent's local search history, and the reassignment phase recombines inter-cluster boundaries.

\begin{figure}[htbp]
\centering
\fbox{\parbox{0.86\linewidth}{\centering Figure 4.2 placeholder. Original image was not present in the txt export.}}
\caption{Illustration of the crossover operator on a nine-station example with $\left|K\right|=4$. Parent ${{\phi }_{A}}$ partitions the stations into $C_{1}^{A},\ldots ,C_{4}^{A},$ with medoids $\{2,4,6,8\}$; parent ${{\phi }_{B}}$ partitions them into $C_{1}^{B},\ldots ,C_{4}^{B}$. The sampled index set is $\mathcal{I}=\{1,3\}$, so the offspring inherits $C_{1}^{A}$ and $C_{3}^{A}$ together with their medoids ${{m}_{1}}=2$ and ${{m}_{3}}=6$. Each remaining station is then allocated from ${{\phi }_{B}}$: stations 5 and 9 lie in clusters whose index is outside $\mathcal{I}$ and are copied directly into ${{C}_{2}}$ and ${{C}_{4}}$ (Line 9 of Algorithm 2); stations 4 and 8 lie in clusters whose index coincides with an inherited one and therefore fall back to the inherited cluster whose medoid minimizes ${{D}_{i,{{m}_{k}}}}$ (Line 11), under the assumption that ${{D}_{4,2}}<{{D}_{4,6}}$ and ${{D}_{8,6}}<{{D}_{8,2}}$.}
\end{figure}

\subsection{Cluster size repair}

Solution construction (Section 4.4) and crossover (Section 4.5) can both produce partitions in which one or more clusters violate the size bounds ${{N}_{\min }}$ and ${{N}_{\max }}$. A two-phase repair procedure restores feasibility before education (Algorithm 1, Line 8).

In Phase 1, the procedure corrects oversized clusters by repeatedly relocating the station most dissimilar to its cluster’s medoid to the nearest cluster whose size is below ${{N}_{\max }}$. In Phase 2, it corrects undersized clusters by pulling the station nearest (by dissimilarity) to the undersized cluster’s medoid from a cluster whose size exceeds ${{N}_{\min }}$. Medoids are recomputed after both phases.

\subsection{Education}

Every new individual undergoes local search education before entering the population (Algorithm 1, Line 9). The procedure applies six operators in a randomly shuffled order and returns as soon as any operator produces an improving move.

Two mechanisms guide candidate selection across all operators. The first is a complementary fit score that estimates a station’s compatibility with a target cluster before committing to an ATS evaluation. For a station $i$ and a target cluster ${{C}_{k}}$, we define the complementary fit score

\[
FS(i,{{C}_{k}})=\frac{1}{|{{{\bar{C}}}_{k}}(i)|}\sum\limits_{j\in {{{\bar{C}}}_{k}}(i)}{(1-{{D}_{ij}})},
\]

where ${{\bar{C}}_{k}}(i)$ denotes the set of stations in ${{C}_{k}}$ whose surplus/deficit status is complementary to that of $i$: surplus stations are scored against deficit members, and vice versa. A higher value indicates stronger UDF complementarity between station $i$ and the members of ${{C}_{k}}$. Balanced stations receive a score of zero and are excluded from candidate generation.

The second mechanism is a granular neighbourhood restriction adapted from Vidal (2022). For each station, a fixed number of nearest neighbours by dissimilarity ${{D}_{ij}}$ are precomputed once before the algorithm begins. LS-1 and LS-2 restrict their search to clusters containing at least one of these nearest neighbours, reducing the set of target clusters from $K$ to a small subset.

Within each operator, candidates are evaluated sequentially by computing the ATS objective for the affected clusters before and after the move, and the first candidate yielding an improvement is accepted. Cluster evaluations are cached by station membership, allowing tentatively rejected moves to still populate the cache for subsequent lookups. The six operators are described below.

LS-1 (Relocate station). Move a single station $i$ from ${{C}_{u}}$ to ${{C}_{v}}$, producing ${{C}_{{{u}'}}}={{C}_{u}}\backslash \{i\}$ and ${{C}_{{{v}'}}}={{C}_{v}}\cup \{i\}$. Granular search restricts candidates to clusters containing at least one of $i$‘s nearest neighbours. A candidate is admitted when $FS(i,{{C}_{v}})>FS(i,{{C}_{u}})$, $|{{C}_{u}}|>{{N}_{\min }}$, and $|{{C}_{v}}|<{{N}_{\max }}$.

LS-2 (Swap stations). Exchange station $i\in {{C}_{u}}$ with station $j\in {{C}_{v}}$, producing ${{C}_{{{u}'}}}=({{C}_{u}}\backslash \{i\})\cup \{j\}$ and ${{C}_{{{v}'}}}=({{C}_{v}}\backslash \{j\})\cup \{i\}$. Granular search applies. A candidate is admitted when $FS(i,{{C}_{v}})>FS(i,{{C}_{u}})$ and $FS(j,{{C}_{u}})>FS(j,{{C}_{v}})$.

LS-3 (Merge and split). Select two clusters ${{C}_{u}}$ and ${{C}_{v}}$ at random, merge their stations into a single pool, and repartition via farthest-pair seeding followed by fit-score-based assignment. This operator restructures two clusters entirely, enabling large-scale partition changes that incremental operators cannot reach. The size bound also applies when admitting a candidate.

LS-4 (Relocate pair). Move station $i\in {{C}_{u}}$ together with its nearest cluster neighbour $j=\arg {{\min }_{{j}'\in {{C}_{u}}\backslash \{i\}}}{{D}_{i{j}'}}$ to ${{C}_{v}}$, producing ${{C}_{{{u}'}}}={{C}_{u}}\backslash \left\{ i,j \right\}$ and ${{C}_{{{v}'}}}={{C}_{v}}\cup \left\{ i,j \right\}$. A candidate is admitted when $FS(i,{{C}_{v}})+FS(j,{{C}_{v}})>FS(i,{{C}_{u}})+FS(j,{{C}_{u}})$, $|{{C}_{u}}|-2\ge {{N}_{\min }}$, and $|{{C}_{v}}|+2\le {{N}_{\max }}$.

LS-5 (Swap pair for single). Exchange a pair $(i,j)$ from ${{C}_{v}}$ with a single station $k$ from ${{C}_{u}}$, producing ${{C}_{{{u}'}}}=({{C}_{u}}\backslash \left\{ k \right\})\cup \left\{ i,j \right\}$ and ${{C}_{{{v}'}}}=({{C}_{v}}\backslash \left\{ i,j \right\})\cup \left\{ k \right\}$. A candidate is admitted when $FS(k,{{C}_{v}})>FS(k,{{C}_{u}})$ and $FS(i,{{C}_{u}})+FS(j,{{C}_{u}})>FS(i,{{C}_{v}})+FS(j,{{C}_{v}})$, with $|{{C}_{u}}|+1\le {{N}_{\max }}$ and $|{{C}_{v}}|-1\ge {{N}_{\min }}$.

LS-6 (Swap pair for pair). Exchange pair $({{i}_{u}},{{j}_{u}})$ from ${{C}_{u}}$ with pair $({{i}_{v}},{{j}_{v}})$ from ${{C}_{v}}$, producing ${{C}_{{{u}'}}}=({{C}_{u}}\backslash \{{{i}_{u}},{{j}_{u}}\})\cup \{{{i}_{v}},{{j}_{v}}\}$ and ${{C}_{{{v}'}}}=({{C}_{v}}\backslash \{{{i}_{v}},{{j}_{v}}\})\cup \{{{i}_{u}},{{j}_{u}}\}$. A candidate is admitted when $FS({{i}_{u}},{{C}_{v}})+FS({{j}_{u}},{{C}_{v}})>FS({{i}_{u}},{{C}_{u}})+FS({{j}_{u}},{{C}_{u}})$ and the symmetric condition holds for $({{i}_{v}},{{j}_{v}})$.

When the fit-score step produces no candidate at all, each operator falls back to evaluating a single random move. This fallback prevents the search from stalling in configurations where the complementary fit criterion excludes all surplus-deficit pairings.

\subsection{Population management}

HGS-UC maintains a single population of feasible solutions whose size ranges between $\mu $ (minimum) and $\mu +\lambda $ (maximum), where $\mu $ is the base population size and $\lambda $ is the number of offspring slots. Selection and survival decisions are driven by a biased fitness measure that balances solution quality and population diversity.

Each individual $\phi$ is assigned a fitness rank $r^{\mathrm{fit}}(\phi)$, determined by objective value in descending order, and a diversity rank $r^{\mathrm{div}}(\phi)$, determined by diversity contribution in descending order. The biased fitness of individual $\phi$ is given by
\[
BF(\phi)=r^{\mathrm{fit}}(\phi)+\left(1-\frac{n_{\mathrm{elite}}}{\mu}\right)r^{\mathrm{div}}(\phi),
\]
where lower values are better.

The diversity contribution of an individual is the average distance to its ${{n}_{\text{close}}}$ nearest neighbours in the population. For two solutions ${{\phi }_{a}}$ and ${{\phi }_{b}}$, let ${{\chi }_{\phi }}\left( i,j \right)=1$ if stations $i$ and $j$ belong to the same cluster in partition $\phi $ and $0$ otherwise. The distance between ${{\phi }_{a}}$ and ${{\phi }_{b}}$ is the fraction of station pairs on which the two partitions disagree,

\[
d\left( {{\phi }_{a}},{{\phi }_{b}} \right)={{\left( \begin{matrix}
\left| V \right|  \\
2  \\
\end{matrix} \right)}^{-1}}\underset{\left\{ i,j \right\}\subseteq V}{\mathop \sum }\,\left| {{\chi }_{{{\phi }_{a}}}}\left( i,j \right)-{{\chi }_{{{\phi }_{b}}}}\left( i,j \right) \right|,
\]

Exact evaluation is $O\left( {{\left| V \right|}^{2}} \right)$ per solution pair, which is computationally expensive given the frequency of distance queries during selection and survival. The implementation uses a sampled estimator $\hat{d}\left( {{\phi }_{a}},{{\phi }_{b}} \right)$ that draws ${{N}_{\text{sample}}}$ station pairs uniformly at random and returns the proportion of broken pairs among them. Each draw is an independent Bernoulli trial with success probability $d({{\phi }_{a}},{{\phi }_{b}})$. By Hoeffding’s inequality (Hoeffding, 1963), we have

\[
\Pr [|\hat{d}-d|\ge \epsilon ]\le 2{{e}^{-2{{N}_{\text{sample}}}{{\epsilon }^{2}}}}\quad \forall \epsilon >0.
\]

Consequently, with ${{N}_{\text{sample}}}=500$ and $\mathsf{\epsilon }=0.1$ used in the implementation, the probability that $\hat{d}$ deviates from $d\left( {{\phi }_{a}},{{\phi }_{b}} \right)$ by more than 0.1 is bounded by $2{{e}^{-10}}<{{10}^{-4}}$. This precision is adequate for the ranking-based use of $d$ in biased fitness.

Parents are selected by binary tournament on biased fitness. Two individuals are drawn at random, and the one with lower biased fitness is chosen. Each generation draws two parents independently. When offspring insertion brings the population size above $\mu +\lambda $, survivor selection removes individuals with the highest biased fitness until the size returns to $\mu $. No explicit clone removal is applied.

The diversity component of biased fitness implicitly penalizes near-duplicate solutions, since individuals surrounded by similar neighbours receive low diversity contributions. Directed diversification is triggered after $I{{t}_{\text{div}}}$ consecutive generations without improvement (Algorithm 1, Lines 21–24). The $\text{max}\left( 1,\left\lfloor \mu /3 \right\rfloor  \right)$ best individuals in the current population are retained by objective value. The remaining population slots are regenerated through the construction procedure of Section 4.4 until the population reaches $4\mu $ individuals. The diversification counter $ite{{r}_{\text{div}}}$ is then reset to zero. The global non-improvement counter $ite{{r}_{\text{NI}}}$ is reset only if diversification yields a new best solution. Otherwise, it continues to accumulate.

\section{Computational experiments}

The HGS-UC algorithm is implemented in single-threaded C++ and compiled with GCC. The clustering MILP, baseline clustering methods, and routing-phase BRP solver are implemented in Python using Gurobi 13.0.1 as the MIP solver. All computational experiments were carried out on a desktop PC equipped with a 13-th Gen Intel Core i9-13900 CPU (24 cores, 32 threads) and 64 GB of RAM. Every Gurobi solve, including the MV-BRP, the clustering MILP, and the per-cluster single-vehicle BRPs, is executed with a single thread to ensure reproducible and comparable timing results across methods. For decomposition-based methods (Gurobi-UC, HGS-UC, and all baseline clustering methods in Section 5.3), the per-cluster single-vehicle BRPs produced by the clustering phase are solved in parallel, each as an independent single-threaded Gurobi process. Accordingly, the reported CPU time of a decomposition-based method is the sum of (i) the clustering time and (ii) the wall-clock time of the longest single-cluster BRP solve within that run, reflecting the end-to-end time an operator would observe. For stochastic methods, we report the mean over 20 independent seeds. The HGS-UC clustering time is the mean over 20 seeds, and the per-run BRP solve time is the maximum across all clusters within that run, averaged over the 20 seeds.

All instances were generated from the New York Citi Bike system data. Citi Bike NYC November 2024 trip records\footnote{\url{https://s3.amazonaws.com/tripdata/index.html}} were aggregated into five-minute pick-up and drop-off rates for each station, and the queueing-theoretic model of Raviv and Kolka (2013) was applied to compute the UDF at every station. After filtering stations with incomplete or anomalous activity records, 1,000 stations were retained as the master pool. To examine how the proposed framework scales with network size, instances of varying sizes $|V|\in \{50,100,200,400,800,1000\}$ were constructed by sampling stations uniformly at random without replacement from the master pool. The duration matrices of the networks were calculated by the OSRM project\footnote{\url{https://project-osrm.org/}}, which provides the duration of the fastest route between all pairs of supplied coordinates. The resulting matrices are asymmetric.

The number of vehicles is $\left| K \right|=\left\lfloor |V|/25 \right\rfloor $, yielding approximately 25 stations per cluster, which is in line with the per-vehicle workload in Schuijbroek et al. (2017). The cluster size bounds are set to ${N_{\min}} = 10$ and ${N_{\max}} = 40$, chosen so that the single-vehicle BRP within each cluster remains tractable for Gurobi at the routing phase. The default vehicle capacity is $Q=20$ and the default repositioning time budget is T = 120 min, unless stated otherwise.

Three sets of experiments were conducted. Section 5.1 describes parameter tuning of HGS-UC on a canonical instance. Section 5.2 compares the proposed CFRS decomposition against the MV-BRP formulation solved directly by Gurobi on small-to-medium instances ($\mid V\mid \in \left\{ 50,100 \right\}$). Section 5.3 evaluates HGS-UC against five baseline clustering methods on large-scale instances ($\mid V\mid \in \left\{ 200,400,800,1000 \right\}$) under varying time budgets and vehicle capacities.

\subsection{Parameter tuning}

The six key HGS-UC parameters were tuned sequentially on a canonical instance with $\mid V\mid =200$ and $\left| K \right|=8$. For each parameter, 20 independent seeds were run per candidate value. The selection criterion was the mean UDF reduction, with mean CPU time as a tiebreaker when the difference in UDF reduction was marginal. Table 5.1 defines the parameters and lists the selected values. Fig. 5.1 shows the mean UDF reduction and CPU time for each candidate value.

\begin{table}[htbp]
\centering
\caption{HGS-UC parameter setting.}
\label{tab:hgsuc-parameters}
\begin{tabular}{p{0.22\linewidth}p{0.56\linewidth}r}
\toprule
Parameter & Description & Value \\
\midrule
$I{{t}_{\text{NI}}}$ & Maximum consecutive iterations without improvement & 1000 \\
$\mu$ & Minimum subpopulation size & 10 \\
$\lambda$ & Number of offspring per generation & 30 \\
${{n}_{\text{elite}}}$ & Number of elite individuals & 6 \\
${{n}_{\text{close}}}$ & Number of close individuals for diversity measurement & 5 \\
$\alpha$ & Dissimilarity trade-off weight & 0.5 \\
\bottomrule
\end{tabular}
\end{table}

\begin{figure}[htbp]
\centering
\fbox{\parbox{0.86\linewidth}{\centering Figure 5.1 placeholder. Original image was not present in the txt export.}}
\caption{Parameter tuning results ($\mid V\mid =200$, $\left| K \right|=8$, 20 seeds per candidate value). Solid lines (left axis): mean UDF reduction. Dashed lines (right axis): mean CPU time.}
\end{figure}

The maximum number of iterations without improvement, $I{{t}_{\text{NI}}}$, exhibits a clear trade-off between solution quality and computation time. The mean UDF reduction increases from 268.06 at $I{{t}_{\text{NI}}}=100$ to 287.60 at $I{{t}_{\text{NI}}}=1000$, after which the objective declines slightly while CPU time continues to grow. The dissimilarity trade-off parameter $\alpha $ shows that a balanced setting ($\alpha =0.5$) outperforms extreme values, confirming that both the travel-time proximity and UDF-improvement potential components contribute to clustering quality. The remaining four parameters ($\mu $, $\lambda $, ${{n}_{\text{elite}}}$, ${{n}_{\text{close}}}$) show relatively small differences across candidates. The selected values are listed in Table 5.1.

\subsection{Three-way benchmark comparison}

Three approaches were compared on instances with $\mid V\mid =50$ ($\left| K \right|=2$) and $\mid V\mid =100$ ($\left| K \right|=4$):

\begin{enumerate}[label=(\roman*)]
\item MV-BRP solved directly by Gurobi with a 2-hour time limit;
\item Gurobi-UC, which solves the exact MILP clustering model (Section 3) at the clustering phase and then solves the per-cluster BRP at the routing phase; and
\item HGS-UC, which replaces the exact clustering with the heuristic described in Section 4, followed by the same routing-phase solver.
\end{enumerate}

Ten instances were generated for each network size. Both $T=120$ min and $T=180$ min were evaluated. HGS-UC was run 20 times per instance with independent seeds. Tables 5.2 and 5.3 present the results. A negative gap indicates that the decomposition method surpasses the MV-BRP incumbent.

\begin{landscape}
\begin{table}[p]
\centering
\caption{Three-way benchmark comparison results, $\mid V\mid =50$ and $\left| K \right|=2$.}
\label{tab:benchmark-50}
\scriptsize
\setlength{\tabcolsep}{3pt}
\begin{adjustbox}{max width=\linewidth}
\begin{tabular}{llrrrrrrr rrrr}
\toprule
\multirow{2}{*}{$T$ (min)} & \multirow{2}{*}{Instance} & \multicolumn{4}{c}{MV-BRP} & \multicolumn{3}{c}{Gurobi-UC} & \multicolumn{4}{c}{HGS-UC} \\
\cmidrule(lr){3-6}\cmidrule(lr){7-9}\cmidrule(lr){10-13}
& & UB & LB & Gap (\%) & CPU (s) & Obj. & Gap* (\%) & CPU (s) & Best & Avg. & Gap** (\%) & CPU (s) \\
\midrule
120 & 1 & 70.76 & 70.76 & 0.00 & 30.60 & 64.86 & 8.34 & 4.80 & 70.76 & 70.36 & 0.57 & 65.40 \\
 & 2 & 66.55 & 66.55 & 0.00 & 608.50 & 65.88 & 1.01 & 2.30 & 66.55 & 66.42 & 0.20 & 84.50 \\
 & 3 & 73.01 & 73.01 & 0.00 & 223.00 & 72.62 & 0.53 & 3.50 & 72.84 & 72.62 & 0.53 & 76.50 \\
 & 4 & 52.61 & 52.61 & 0.00 & 117.90 & 52.11 & 0.95 & 3.20 & 52.26 & 52.11 & 0.95 & 85.70 \\
 & 5 & 60.84 & 60.84 & 0.00 & 2489.40 & 59.57 & 2.09 & 7.00 & 60.82 & 60.50 & 0.56 & 24.80 \\
 & 6 & 37.22 & 37.22 & 0.00 & 648.40 & 36.10 & 3.01 & 2.00 & 37.09 & 36.90 & 0.86 & 10.90 \\
 & 7 & 57.05 & 57.05 & 0.00 & 355.20 & 56.15 & 1.58 & 2.20 & 56.88 & 56.74 & 0.54 & 8.80 \\
 & 8 & 43.11 & 43.11 & 0.00 & 3075.70 & 42.30 & 1.88 & 3.60 & 43.11 & 43.05 & 0.14 & 11.30 \\
 & 9 & 48.01 & 48.01 & 0.00 & 636.40 & 47.21 & 1.67 & 3.10 & 48.01 & 47.82 & 0.40 & 9.00 \\
 & 10 & 49.13 & 49.13 & 0.00 & 199.40 & 48.56 & 1.16 & 2.90 & 49.08 & 48.82 & 0.63 & 9.00 \\
\midrule
 & Avg &  &  & 0.00 & 838.45 &  & 2.22 & 3.46 &  &  & 0.54 & 38.59 \\
180 & 1 & 74.82 & 74.82 & 0.00 & 497.00 & 73.62 & 1.60 & 9.00 & 74.82 & 74.44 & 0.51 & 53.30 \\
 & 2 & 69.17 & 69.17 & 0.00 & 862.10 & 68.20 & 1.40 & 3.70 & 68.89 & 68.79 & 0.55 & 131.60 \\
 & 3 & 76.56 & 76.56 & 0.00 & 1313.90 & 75.82 & 0.97 & 7.20 & 76.28 & 75.97 & 0.77 & 58.80 \\
 & 4 & 63.21 & 63.12 & 0.14 & 7200.00 & 60.67 & 3.88 & 3.20 & 63.12 & 63.03 & 0.14 & 63.90 \\
 & 5 & 74.93 & 74.69 & 0.31 & 7200.00 & 73.39 & 1.74 & 9.20 & 74.60 & 74.33 & 0.48 & 18.10 \\
 & 6 & 38.52 & 38.52 & 0.00 & 970.10 & 38.20 & 0.83 & 5.60 & 38.52 & 38.31 & 0.55 & 59.50 \\
 & 7 & 60.49 & 60.49 & 0.00 & 2156.80 & 59.84 & 1.07 & 7.00 & 60.49 & 60.40 & 0.15 & 60.30 \\
 & 8 & 46.76 & 46.66 & 0.21 & 7200.00 & 46.44 & 0.47 & 5.30 & 46.66 & 46.39 & 0.58 & 55.00 \\
 & 9 & 51.41 & 51.41 & 0.00 & 3446.20 & 51.22 & 0.36 & 6.10 & 51.41 & 51.22 & 0.37 & 22.00 \\
 & 10 & 50.91 & 50.91 & 0.00 & 852.10 & 50.51 & 0.78 & 11.40 & 50.76 & 50.56 & 0.69 & 35.40 \\
\midrule
 & Avg &  &  & 0.07 & 3169.82 &  & 1.31 & 6.77 &  &  & 0.48 & 55.79 \\
\bottomrule
\end{tabular}
\end{adjustbox}
\begin{tablenotes}\footnotesize
\item *Gap = (MV-BRP LB -- Gurobi-UC Obj.) / MV-BRP LB $\times$ 100.
\item **Gap = (MV-BRP LB -- HGS-UC Avg.) / MV-BRP LB $\times$ 100.
\end{tablenotes}
\end{table}
\end{landscape}

\begin{landscape}
\begin{table}[p]
\centering
\caption{Three-way benchmark comparison results, $\mid V\mid =100$ and $\left| K \right|=4$.}
\label{tab:benchmark-100}
\scriptsize
\setlength{\tabcolsep}{3pt}
\begin{adjustbox}{max width=\linewidth}
\begin{tabular}{llrrrrrrr rrrr}
\toprule
\multirow{2}{*}{$T$ (min)} & \multirow{2}{*}{Instance} & \multicolumn{4}{c}{MV-BRP} & \multicolumn{3}{c}{Gurobi-UC} & \multicolumn{4}{c}{HGS-UC} \\
\cmidrule(lr){3-6}\cmidrule(lr){7-9}\cmidrule(lr){10-13}
& & UB & LB & Gap (\%) & CPU (s) & Obj. & Gap* (\%) & CPU (s) & Best & Avg. & Gap** (\%) & CPU (s) \\
\midrule
120 & 1 & 91.11 & 86.24 & 5.35 & 7200.00 & 86.50 & -0.30 & 1328.80 & 87.08 & 86.89 & -0.75 & 36.70 \\
 & 2 & 107.55 & 98.38 & 8.53 & 7200.00 & 96.54 & 1.87 & 2202.30 & 100.73 & 100.37 & -2.02 & 53.90 \\
 & 3 & 89.37 & 83.57 & 6.50 & 7200.00 & 82.82 & 0.90 & 7210.50 & 83.18 & 83.08 & 0.59 & 253.20 \\
 & 4 & 137.02 & 122.08 & 10.90 & 7200.00 & 124.65 & -2.10 & 7202.60 & 125.56 & 125.29 & -2.63 & 51.30 \\
 & 5 & 173.62 & 157.64 & 9.20 & 7200.00 & 152.99 & 2.95 & 7203.00 & 160.92 & 160.46 & -1.79 & 42.50 \\
 & 6 & 139.68 & 128.78 & 7.81 & 7200.00 & 127.91 & 0.67 & 7203.10 & 129.64 & 129.23 & -0.35 & 86.30 \\
 & 7 & 75.72 & 71.98 & 4.94 & 7200.00 & 71.28 & 0.97 & 7204.00 & 72.16 & 72.05 & -0.10 & 74.80 \\
 & 8 & 157.71 & 138.75 & 12.02 & 7200.00 & 140.18 & -1.03 & 7206.90 & 141.65 & 141.33 & -1.86 & 51.90 \\
 & 9 & 131.62 & 127.57 & 3.07 & 7200.00 & 123.82 & 2.94 & 5198.60 & 127.62 & 127.13 & 0.34 & 38.40 \\
 & 10 & 112.77 & 104.60 & 7.24 & 7200.00 & 101.88 & 2.61 & 7204.70 & 105.88 & 105.62 & -0.98 & 52.40 \\
\midrule
 & Avg &  &  & 7.56 & 7200.00 &  & 0.95 & 5916.45 &  &  & -0.95 & 74.14 \\
180 & 1 & 92.20 & 90.34 & 2.02 & 7200.00 & 90.86 & -0.58 & 7232.70 & 91.85 & 91.66 & -1.46 & 264.10 \\
 & 2 & 110.68 & 108.70 & 1.79 & 7200.00 & 108.64 & 0.06 & 7228.00 & 109.63 & 109.25 & -0.51 & 77.00 \\
 & 3 & 89.40 & 87.77 & 1.82 & 7200.00 & 88.06 & -0.33 & 7222.30 & 89.32 & 88.95 & -1.34 & 108.10 \\
 & 4 & 143.41 & 136.73 & 4.66 & 7200.00 & 140.57 & -2.81 & 7224.60 & 141.52 & 141.24 & -3.30 & 158.80 \\
 & 5 & 194.97 & 186.42 & 4.38 & 7200.00 & 182.76 & 1.96 & 7222.20 & 188.13 & 187.71 & -0.69 & 228.90 \\
 & 6 & 158.83 & 151.09 & 4.88 & 7200.00 & 150.13 & 0.63 & 7215.10 & 154.45 & 154.09 & -1.99 & 126.10 \\
 & 7 & 76.28 & 74.74 & 2.02 & 7200.00 & 75.08 & -0.44 & 7232.40 & 76.28 & 75.87 & -1.51 & 160.20 \\
 & 8 & 179.13 & 170.69 & 4.71 & 7200.00 & 165.87 & 2.83 & 7229.00 & 170.93 & 170.77 & -0.05 & 148.40 \\
 & 9 & 135.84 & 131.91 & 2.89 & 7200.00 & 133.65 & -1.31 & 7205.10 & 135.32 & 134.87 & -2.24 & 68.30 \\
 & 10 & 118.78 & 114.97 & 3.21 & 7200.00 & 116.00 & -0.90 & 7231.00 & 117.53 & 117.20 & -1.94 & 81.40 \\
\midrule
 & Avg &  &  & 3.24 & 7200.00 &  & -0.09 & 7224.24 &  &  & -1.50 & 142.13 \\
\bottomrule
\end{tabular}
\end{adjustbox}
\begin{tablenotes}\footnotesize
\item *Gap = (MV-BRP LB -- Gurobi-UC Obj.) / MV-BRP LB $\times$ 100.
\item **Gap = (MV-BRP LB -- HGS-UC Avg.) / MV-BRP LB $\times$ 100.
\end{tablenotes}
\end{table}
\end{landscape}

For the MV-BRP, UB and LB denote the upper bound (relaxation bound) and lower bound (best feasible solution) returned by Gurobi, respectively, and the gap is the Gurobi optimality gap (UB − LB) / LB × 100. For the decomposition methods, Best and Avg. denote the maximum and mean UDF reduction across seeds (for the deterministic Gurobi-UC, Best and Avg. coincide with its single objective value). The gap is computed as (MV-BRP LB − Method Obj.) / MV-BRP LB × 100, where Method Obj. is the Gurobi-UC objective or HGS-UC average, as indicated in the table footnotes. A negative value indicates that the method surpasses the MV-BRP incumbent.

At $\mid V\mid =50$, the MV-BRP solved all instances to optimality within the 2-hour time limit at $T=120$ min and most instances at $T=180$ min (average Gurobi gap 0.07\%). Gurobi-UC solved the clustering MILP within seconds and achieved an average gap of 2.22\% at $T=120$ min and 1.31\% at $T=180$ min. HGS-UC achieved an average gap of 0.54\% and 0.48\%, respectively, with the best seed matching the MV-BRP incumbent solution on several instances. At $\mid V\mid =100$, the MV-BRP hit the 2-hour time limit on all 10 instances at both time budgets, with an average Gurobi gap of 7.56\% at $T=120$ min and 3.24\% at $T=180$ min. HGS-UC surpassed the MV-BRP incumbent on 8 out of 10 instances at $T=120$ min (average gap $-0.95\%$, average CPU time 74.1 s) and on all 10 instances at $T=180$ min (average gap -1.50\%, average CPU time 142.1 s). Gurobi-UC also produced competitive solutions at this scale (average gap 0.95\% and -0.09\% at the two time budgets), but its total CPU time approached or exceeded the MV-BRP time limit because the clustering MILP itself becomes intractable at $\mid V\mid =100$. These results confirm that the CFRS decomposition is a viable strategy for the multi-vehicle BRP, and that HGS-UC provides both competitive solution quality and a substantial reduction in computation time.

\subsection{Baseline comparison}

For each of $\mid V\mid \in \left\{ 200,400,800,1000 \right\}$, one instance was generated and each stochastic method was run 20 times with independent seeds. Five baseline clustering methods were compared against HGS-UC:

\begin{itemize}
\item K-means: Multidimensional scaling is applied to the travel-time matrix, followed by standard K-means clustering in the resulting coordinate space.
\item KNN: A $k$-nearest-neighbour graph is constructed from travel times and spectral clustering is applied.
\item Schuijbroek et al. (2017): Stations are clustered to minimize the total service-level violation, accounting for inventory imbalance at each station.
\item Forma et al. (2015): A savings-based heuristic partitions stations into vehicle districts by merging station pairs that share high UDF-reduction potential.
\item Lv et al. (2020): A self-sufficiency-based clustering method that constructs clusters where internal supply and demand are balanced. The perturbation parameter $\theta $ was not specified in the original paper. We tested $\theta \in \left\{ 0,5,10,20 \right\}$.
\end{itemize}

All methods use the same routing-phase BRP solver and the same instance data. The cluster size is $\left| K \right|=\left\lfloor \mid V\mid /25 \right\rfloor $ for all methods. Table 5.4 reports the average UDF reduction (Avg.), the best single run (Best), the CPU time, and the percentage gap relative to HGS-UC, defined as (Method Avg. – HGS-UC Avg.) / HGS-UC Avg. × 100. A negative gap indicates that HGS-UC outperforms the baseline. The reported performance therefore reflects the cluster compositions produced by each method under a common routing benchmark.

\begin{landscape}
\begin{longtable}{llp{4.3cm}rrrrr}
\caption{Baseline comparison results ($T=120$ min, $Q=20$, 20 seeds per stochastic method).}\label{tab:baseline-default}\\
\toprule
$|V|$ & $|K|$ & Method & Avg. & Best & CPU (s) & Gap* (\%) & p-value** \\
\midrule
\endfirsthead
\toprule
$|V|$ & $|K|$ & Method & Avg. & Best & CPU (s) & Gap* (\%) & p-value** \\
\midrule
\endhead
\bottomrule
\multicolumn{8}{l}{\footnotesize *Gap = (Method Avg. -- HGS-UC Avg.) / HGS-UC Avg. $\times$ 100.}\\
\multicolumn{8}{l}{\footnotesize **p-values obtained from one-sided Welch's t-test.}\\
\endfoot
200 & 8 & HGS-UC & 286.23 & 289.62 & 324.33 & -- & -- \\
 &  & Schuijbroek et al. & 244.06 & 244.06 & 65.82 & -14.73 & 1.90E-23 \\
 &  & Forma et al. & 232.42 & 232.42 & 5.86 & -18.80 & 1.89E-25 \\
 &  & Lv et al. ($\theta$ = 0) & 207.56 & 209.93 & 31.64 & -27.48 & 1.12E-25 \\
 &  & Lv et al. ($\theta$ = 5) & 205.81 & 222.26 & 40.77 & -28.09 & 1.40E-16 \\
 &  & Lv et al. ($\theta$ = 10) & 198.44 & 219.63 & 47.15 & -30.67 & 7.33E-15 \\
 &  & Lv et al. ($\theta$ = 20) & 200.87 & 220.47 & 54.68 & -29.82 & 2.33E-13 \\
 &  & K-means & 212.81 & 215.53 & 11.25 & -25.65 & 1.97E-24 \\
 &  & KNN & 211.51 & 224.26 & 29.63 & -26.10 & 1.36E-26 \\
\midrule
400 & 16 & HGS-UC & 494.29 & 508.01 & 773.81 & -- & -- \\
 &  & Schuijbroek et al. & 406.77 & 406.77 & 89.71 & -17.71 & 6.66E-14 \\
 &  & Forma et al. & 381.82 & 381.82 & 19.73 & -22.75 & 1.39E-24 \\
 &  & Lv et al. ($\theta$ = 0) & 293.26 & 333.00 & 46.46 & -40.67 & 2.68E-19 \\
 &  & Lv et al. ($\theta$ = 5) & 307.72 & 356.14 & 64.54 & -37.74 & 1.86E-18 \\
 &  & Lv et al. ($\theta$ = 10) & 301.83 & 339.41 & 72.67 & -38.94 & 1.88E-18 \\
 &  & Lv et al. ($\theta$ = 20) & 297.06 & 333.04 & 87.92 & -39.90 & 5.14E-21 \\
 &  & K-means & 315.75 & 354.08 & 14.05 & -36.12 & 2.68E-18 \\
 &  & KNN & 318.89 & 319.09 & 29.37 & -35.48 & 2.23E-28 \\
\midrule
800 & 32 & HGS-UC & 1068.71 & 1085.21 & 2227.90 & -- & -- \\
 &  & Schuijbroek et al. & 884.36 & 884.36 & 313.25 & -17.25 & 1.62E-28 \\
 &  & Forma et al. & 865.84 & 865.84 & 20.09 & -18.98 & 5.51E-28 \\
 &  & Lv et al. ($\theta$ = 0) & 709.11 & 767.39 & 80.64 & -33.65 & 1.11E-19 \\
 &  & Lv et al. ($\theta$ = 5) & 736.28 & 773.94 & 113.47 & -31.10 & 9.70E-22 \\
 &  & Lv et al. ($\theta$ = 10) & 711.75 & 773.15 & 132.54 & -33.40 & 1.34E-19 \\
 &  & Lv et al. ($\theta$ = 20) & 700.01 & 752.04 & 162.33 & -34.50 & 1.19E-22 \\
 &  & K-means & 697.73 & 733.97 & 27.93 & -34.71 & 3.74E-23 \\
 &  & KNN & 700.35 & 735.59 & 29.21 & -34.47 & 4.01E-27 \\
\midrule
1000 & 40 & HGS-UC & 1429.16 & 1453.92 & 3228.15 & -- & -- \\
 &  & Schuijbroek et al. & 1160.66 & 1160.66 & 554.44 & -18.79 & 1.92E-26 \\
 &  & Forma et al. & 961.75 & 961.75 & 27.19 & -32.71 & 5.18E-31 \\
 &  & Lv et al. ($\theta$ = 0) & 893.02 & 954.43 & 98.18 & -37.51 & 2.01E-24 \\
 &  & Lv et al. ($\theta$ = 5) & 913.83 & 986.05 & 135.78 & -36.06 & 7.98E-23 \\
 &  & Lv et al. ($\theta$ = 10) & 893.76 & 961.88 & 159.03 & -37.46 & 4.55E-26 \\
 &  & Lv et al. ($\theta$ = 20) & 890.52 & 952.22 & 201.56 & -37.69 & 5.05E-24 \\
 &  & K-means & 874.83 & 941.21 & 41.82 & -38.79 & 5.80E-23 \\
 &  & KNN & 847.97 & 873.94 & 28.62 & -40.67 & 7.13E-31 \\
\end{longtable}
\end{landscape}

HGS-UC achieves the highest average UDF reduction at every instance size. The gap ranges from -14.73\% (Schuijbroek et al., $\mid V\mid =200$) to -40.67\% (KNN, $\mid V\mid =1000$). All differences are statistically significant (one-sided Welch’s t-test, $p<{{10}^{-12}}$ in all cases). The HGS-UC clustering requires more CPU time than the baseline methods because it solves an optimization-based clustering problem rather than applying a closed-form partitioning rule. The total CPU time of the CFRS pipeline (clustering plus routing) is approximately 54 minutes, which fits within the overnight repositioning window.

Tables 5.5 and 5.6 report the gap of each baseline method relative to HGS-UC under varying time budgets ($T\in \left\{ 60,120,180,240 \right\}$ min with $Q=20$) and vehicle capacities ($Q\in \left\{ 10,20,30,40 \right\}$ with $T=120$ min), respectively. Methods whose clustering depends on T or Q (HGS-UC, Schuijbroek et al., and Lv et al. for T; Schuijbroek et al. and Lv et al. for Q) were re-clustered at each parameter value. HGS-UC clustering is independent of Q and was not re-clustered for the capacity experiments.

\begin{landscape}
\begin{longtable}{p{4.2cm}lrrrrrrrr}
\caption{Baseline comparison results under varying time budgets ($Q=20$, 20 seeds per stochastic method).}\label{tab:baseline-time}\\
\toprule
\multirow{2}{*}{Method} & \multirow{2}{*}{$|V|$} & \multicolumn{2}{c}{$T=60$} & \multicolumn{2}{c}{$T=120$} & \multicolumn{2}{c}{$T=180$} & \multicolumn{2}{c}{$T=240$} \\
\cmidrule(lr){3-4}\cmidrule(lr){5-6}\cmidrule(lr){7-8}\cmidrule(lr){9-10}
& & Gap* (\%) & p-value** & Gap (\%) & p-value & Gap (\%) & p-value & Gap (\%) & p-value \\
\midrule
\endfirsthead
\toprule
\multirow{2}{*}{Method} & \multirow{2}{*}{$|V|$} & \multicolumn{2}{c}{$T=60$} & \multicolumn{2}{c}{$T=120$} & \multicolumn{2}{c}{$T=180$} & \multicolumn{2}{c}{$T=240$} \\
\cmidrule(lr){3-4}\cmidrule(lr){5-6}\cmidrule(lr){7-8}\cmidrule(lr){9-10}
& & Gap* (\%) & p-value** & Gap (\%) & p-value & Gap (\%) & p-value & Gap (\%) & p-value \\
\midrule
\endhead
\bottomrule
\multicolumn{10}{l}{\footnotesize *Gap = (Method Avg. -- HGS-UC Avg.) / HGS-UC Avg. $\times$ 100.}\\
\multicolumn{10}{l}{\footnotesize **p-values obtained from one-sided Welch's t-test.}\\
\endfoot
Schuijbroek et al. & 200 & -21.45 & 1.73E-24 & -14.73 & 1.90E-23 & -12.82 & 3.73E-30 & -11.37 & 1.01E-42 \\
 & 400 & -18.44 & 4.48E-07 & -17.71 & 6.66E-14 & -14.82 & 1.11E-18 & -13.19 & 1.69E-25 \\
 & 800 & -18.78 & 1.13E-19 & -17.25 & 1.62E-28 & -16.53 & 6.33E-27 & -13.20 & 7.47E-31 \\
 & 1000 & -19.96 & 1.26E-18 & -18.79 & 1.92E-26 & -18.32 & 1.97E-30 & -16.23 & 9.15E-39 \\
\midrule
Forma et al. & 200 & -31.16 & 1.46E-27 & -18.80 & 1.89E-25 & -7.50 & 9.68E-26 & -2.27 & 1.98E-29 \\
 & 400 & -30.66 & 2.90E-24 & -22.75 & 1.39E-24 & -13.84 & 4.24E-32 & -6.24 & 2.25E-29 \\
 & 800 & -19.16 & 7.71E-20 & -18.98 & 5.51E-28 & -11.60 & 5.25E-24 & -5.00 & 7.65E-23 \\
 & 1000 & -34.25 & 2.07E-25 & -32.71 & 5.18E-31 & -23.27 & 2.11E-32 & -14.33 & 9.81E-38 \\
\midrule
Lv et al. ($\theta$ = 0) & 200 & -33.50 & 2.23E-24 & -27.48 & 1.12E-25 & -19.99 & 1.99E-24 & -12.27 & 8.09E-20 \\
 & 400 & -45.74 & 2.19E-20 & -40.67 & 2.68E-19 & -31.28 & 1.40E-16 & -21.73 & 6.86E-13 \\
 & 800 & -36.63 & 4.73E-20 & -33.65 & 1.11E-19 & -26.29 & 6.70E-17 & -18.23 & 5.15E-15 \\
 & 1000 & -41.70 & 5.03E-24 & -37.51 & 2.01E-24 & -28.83 & 1.90E-20 & -18.89 & 1.64E-16 \\
\midrule
Lv et al. ($\theta$ = 5) & 200 & -34.47 & 9.78E-19 & -28.09 & 1.40E-16 & -18.98 & 8.04E-13 & -9.66 & 3.24E-07 \\
 & 400 & -43.17 & 9.41E-20 & -37.74 & 1.86E-18 & -28.51 & 2.16E-15 & -19.26 & 4.83E-12 \\
 & 800 & -34.96 & 4.07E-21 & -31.10 & 9.70E-22 & -23.75 & 2.18E-20 & -15.61 & 4.90E-16 \\
 & 1000 & -40.39 & 9.69E-24 & -36.06 & 7.98E-23 & -27.53 & 2.05E-20 & -18.01 & 1.04E-16 \\
\midrule
Lv et al. ($\theta$ = 10) & 200 & -35.87 & 4.55E-18 & -30.67 & 7.33E-15 & -22.37 & 5.74E-12 & -13.89 & 2.59E-08 \\
 & 400 & -43.69 & 6.84E-22 & -38.94 & 1.88E-18 & -30.23 & 7.27E-16 & -21.53 & 2.93E-12 \\
 & 800 & -37.52 & 1.20E-19 & -33.40 & 1.34E-19 & -26.02 & 2.91E-17 & -17.33 & 3.58E-14 \\
 & 1000 & -41.87 & 2.65E-26 & -37.46 & 4.55E-26 & -28.86 & 1.48E-20 & -18.89 & 2.58E-16 \\
\midrule
Lv et al. ($\theta$ = 20) & 200 & -35.83 & 8.55E-17 & -29.82 & 2.33E-13 & -20.87 & 2.01E-09 & -12.09 & 1.31E-05 \\
 & 400 & -45.89 & 1.13E-22 & -39.90 & 5.14E-21 & -29.81 & 2.77E-17 & -19.96 & 6.83E-13 \\
 & 800 & -38.41 & 1.04E-20 & -34.50 & 1.19E-22 & -26.75 & 1.47E-22 & -17.69 & 2.87E-18 \\
 & 1000 & -42.07 & 2.81E-24 & -37.69 & 5.05E-24 & -29.07 & 1.64E-19 & -18.99 & 3.69E-15 \\
\midrule
K-means & 200 & -33.61 & 1.04E-26 & -25.65 & 1.97E-24 & -14.37 & 3.27E-23 & -4.24 & 3.00E-22 \\
 & 400 & -42.77 & 6.01E-21 & -36.12 & 2.68E-18 & -27.32 & 3.98E-16 & -18.26 & 3.48E-13 \\
 & 800 & -38.19 & 1.31E-21 & -34.71 & 3.74E-23 & -27.56 & 4.00E-21 & -20.02 & 3.80E-19 \\
 & 1000 & -42.48 & 2.08E-21 & -38.79 & 5.80E-23 & -30.95 & 7.38E-21 & -21.90 & 5.53E-17 \\
\midrule
KNN & 200 & -36.56 & 3.87E-28 & -26.10 & 1.36E-26 & -14.10 & 1.19E-24 & -4.18 & 5.40E-22 \\
 & 400 & -40.60 & 1.37E-26 & -35.48 & 2.23E-28 & -28.29 & 3.32E-38 & -20.72 & 2.61E-39 \\
 & 800 & -39.91 & 1.59E-23 & -34.47 & 4.01E-27 & -26.21 & 5.75E-24 & -17.67 & 2.57E-21 \\
 & 1000 & -45.05 & 6.91E-26 & -40.67 & 7.13E-31 & -33.07 & 2.91E-26 & -24.71 & 1.17E-24 \\
\end{longtable}
\end{landscape}

\begin{landscape}
\begin{longtable}{p{4.2cm}lrrrrrrrr}
\caption{Baseline comparison results under varying vehicle capacities ($T=120$ min, 20 seeds per stochastic method).}\label{tab:baseline-capacity}\\
\toprule
\multirow{2}{*}{Method} & \multirow{2}{*}{$|V|$} & \multicolumn{2}{c}{$Q=10$} & \multicolumn{2}{c}{$Q=20$} & \multicolumn{2}{c}{$Q=30$} & \multicolumn{2}{c}{$Q=40$} \\
\cmidrule(lr){3-4}\cmidrule(lr){5-6}\cmidrule(lr){7-8}\cmidrule(lr){9-10}
& & Gap* (\%) & p-value** & Gap (\%) & p-value & Gap (\%) & p-value & Gap (\%) & p-value \\
\midrule
\endfirsthead
\toprule
\multirow{2}{*}{Method} & \multirow{2}{*}{$|V|$} & \multicolumn{2}{c}{$Q=10$} & \multicolumn{2}{c}{$Q=20$} & \multicolumn{2}{c}{$Q=30$} & \multicolumn{2}{c}{$Q=40$} \\
\cmidrule(lr){3-4}\cmidrule(lr){5-6}\cmidrule(lr){7-8}\cmidrule(lr){9-10}
& & Gap* (\%) & p-value** & Gap (\%) & p-value & Gap (\%) & p-value & Gap (\%) & p-value \\
\midrule
\endhead
\bottomrule
\multicolumn{10}{l}{\footnotesize *Gap = (Method Avg. -- HGS-UC Avg.) / HGS-UC Avg. $\times$ 100.}\\
\multicolumn{10}{l}{\footnotesize **p-values obtained from one-sided Welch's t-test.}\\
\endfoot
Schuijbroek et al. & 200 & -9.60 & 6.28E-17 & -14.73 & 1.90E-23 & -17.31 & 1.82E-17 & -15.02 & 7.59E-21 \\
 & 400 & -14.09 & 3.32E-24 & -17.71 & 6.66E-14 & -15.44 & 1.41E-20 & -8.73 & 1.54E-16 \\
 & 800 & -22.40 & 2.06E-27 & -17.25 & 1.62E-28 & -17.65 & 7.29E-24 & -16.88 & 2.94E-22 \\
 & 1000 & -18.38 & 5.04E-26 & -18.79 & 1.92E-26 & -12.85 & 2.04E-21 & -14.23 & 5.69E-21 \\
\midrule
Forma et al. & 200 & -13.82 & 6.95E-20 & -18.80 & 1.89E-25 & -20.22 & 1.62E-24 & -20.29 & 2.61E-23 \\
 & 400 & -19.80 & 5.25E-27 & -22.75 & 1.39E-24 & -23.22 & 6.35E-24 & -23.32 & 1.48E-24 \\
 & 800 & -18.68 & 3.19E-26 & -18.98 & 5.51E-28 & -18.54 & 2.86E-24 & -18.74 & 4.07E-23 \\
 & 1000 & -30.80 & 2.81E-30 & -32.71 & 5.18E-31 & -33.25 & 3.12E-29 & -32.75 & 7.98E-28 \\
\midrule
Lv et al. ($\theta$ = 0) & 200 & -13.34 & 3.45E-15 & -27.48 & 1.12E-25 & -30.62 & 6.74E-15 & -29.73 & 5.09E-16 \\
 & 400 & -35.97 & 3.70E-20 & -40.67 & 2.68E-19 & -33.91 & 2.11E-16 & -25.91 & 3.34E-17 \\
 & 800 & -28.36 & 2.25E-18 & -33.65 & 1.11E-19 & -30.53 & 1.99E-21 & -25.61 & 1.63E-15 \\
 & 1000 & -33.82 & 3.04E-20 & -37.51 & 2.01E-24 & -34.52 & 7.18E-22 & -25.28 & 5.23E-17 \\
\midrule
Lv et al. ($\theta$ = 5) & 200 & -19.21 & 1.38E-10 & -28.09 & 1.40E-16 & -32.15 & 6.60E-15 & -26.24 & 1.49E-12 \\
 & 400 & -35.49 & 2.10E-20 & -37.74 & 1.86E-18 & -34.54 & 7.84E-18 & -25.80 & 3.59E-15 \\
 & 800 & -26.82 & 1.07E-18 & -31.10 & 9.70E-22 & -27.59 & 4.97E-22 & -24.37 & 1.20E-16 \\
 & 1000 & -32.81 & 3.75E-20 & -36.06 & 7.98E-23 & -33.36 & 1.54E-20 & -23.18 & 2.29E-15 \\
\midrule
Lv et al. ($\theta$ = 10) & 200 & -17.43 & 8.83E-12 & -30.67 & 7.33E-15 & -32.91 & 2.62E-15 & -26.30 & 3.17E-11 \\
 & 400 & -34.73 & 9.00E-18 & -38.94 & 1.88E-18 & -35.12 & 4.02E-18 & -25.31 & 1.78E-16 \\
 & 800 & -26.96 & 2.06E-20 & -33.40 & 1.34E-19 & -29.87 & 1.20E-18 & -24.23 & 1.87E-17 \\
 & 1000 & -33.85 & 9.47E-21 & -37.46 & 4.55E-26 & -32.31 & 1.29E-21 & -21.56 & 3.70E-16 \\
\midrule
Lv et al. ($\theta$ = 20) & 200 & -16.00 & 4.65E-10 & -29.82 & 2.33E-13 & -31.86 & 5.84E-18 & -30.41 & 9.94E-14 \\
 & 400 & -37.56 & 7.26E-18 & -39.90 & 5.14E-21 & -35.19 & 2.03E-17 & -24.66 & 2.29E-13 \\
 & 800 & -28.42 & 5.34E-17 & -34.50 & 1.19E-22 & -30.11 & 2.78E-21 & -20.30 & 9.81E-16 \\
 & 1000 & -33.43 & 3.15E-20 & -37.69 & 5.05E-24 & -33.78 & 7.85E-22 & -23.34 & 5.93E-19 \\
\midrule
K-means & 200 & -13.02 & 3.01E-16 & -25.65 & 1.97E-24 & -28.16 & 8.81E-26 & -29.16 & 1.37E-25 \\
 & 400 & -28.21 & 1.53E-16 & -36.12 & 2.68E-18 & -37.56 & 1.35E-18 & -37.78 & 6.47E-19 \\
 & 800 & -27.16 & 1.25E-20 & -34.71 & 3.74E-23 & -36.14 & 3.13E-23 & -36.66 & 2.63E-23 \\
 & 1000 & -29.79 & 4.96E-22 & -38.79 & 5.80E-23 & -40.46 & 6.75E-23 & -40.95 & 8.24E-23 \\
\midrule
KNN & 200 & -13.93 & 1.44E-19 & -26.10 & 1.36E-26 & -28.62 & 2.55E-27 & -29.60 & 4.78E-27 \\
 & 400 & -28.77 & 2.71E-30 & -35.48 & 2.23E-28 & -36.76 & 8.12E-28 & -36.92 & 1.87E-28 \\
 & 800 & -26.53 & 2.71E-24 & -34.47 & 4.01E-27 & -35.82 & 1.18E-25 & -36.40 & 3.88E-25 \\
 & 1000 & -31.87 & 2.66E-27 & -40.67 & 7.13E-31 & -42.29 & 1.23E-30 & -42.74 & 1.70E-30 \\
\end{longtable}
\end{landscape}

HGS-UC outperforms all baselines at every time-budget level, every capacity level, and every instance size. All gaps are statistically significant ($p<10^{-4}$).

\section{Conclusion}

This paper addresses the multi-vehicle static BRP with UDF as the service measure and proposes a CFRS framework in which the clustering phase is formulated as an MILP that jointly captures UDF reduction and vehicle routing time feasibility. To the best of our knowledge, no prior CFRS framework for the BRP formulates clustering as an optimization model that simultaneously evaluates UDF reduction at station-level inventories and enforces a per-cluster routing time budget. Because this MILP is intractable for city-scale instances, we develop HGS-UC, a hybrid genetic search tailored to the clustering model. HGS-UC represents a solution as a partition of the station set, constructs initial solutions by K-medoids clustering on a dissimilarity measure that combines travel-time proximity and bilateral UDF reduction potential, and refines them through a local search with six operators guided by a complementary fit score and a granular neighborhood restriction. Solution evaluation uses the ATS procedure, which estimates a time-feasible transfer set for each cluster by selecting surplus-deficit transfers in decreasing order of UDF reduction per unit handling time. ATS replaces an exact single-vehicle BRP solve that would be prohibitively expensive at every evaluation, making population-based search tractable on city-scale instances.

Computational experiments on instances generated from New York Citi Bike data demonstrate the framework's performance across three instance scales. First, on the 50-station instances, HGS-UC obtains solutions whose mean is within 0.54\% of the MV-BRP incumbent in under a minute of CPU time, with the best seed matching the incumbent on several instances. Second, on the 100-station instances, where Gurobi reaches the two-hour limit on every instance with average optimality gaps of 7.56\% and 3.24\% at the 120-minute and 180-minute budgets, respectively,  HGS-UC's mean solution exceeds the Gurobi incumbent on 8 of 10 instances at the 120-minute repositioning budget and on all 10 instances at the 180-minute budget, in 1 to 3 minutes of CPU time. Third, on instances of 200 to 1,000 stations — a scale not previously reported in the CFRS literature for the BRP — HGS-UC delivers higher UDF reduction than five baseline clustering methods, namely the methods of Schuijbroek et al. (2017), Forma et al. (2015), and Lv et al. (2020), together with K-means and KNN, at every combination of network size, time budget, and vehicle capacity tested, with all gaps statistically significant ($p<10^{-4}$). At the default time budget and vehicle capacity, the average UDF reduction gains range from 14.73\% to 40.67\%, with consistent improvements across variations in both parameters. On the largest 1,000-station instances, the complete CFRS pipeline (clustering plus per-cluster routing) finishes in under one hour, well within the overnight repositioning window.

\section*{References}
\addcontentsline{toc}{section}{References}
\begingroup
\setlength{\parindent}{0pt}
\setlength{\parskip}{0.45em}
Arthur, D., Vassilvitskii, S., 2007. k-means++: the advantages of careful seeding. In: Proceedings of the Eighteenth Annual ACM-SIAM Symposium on Discrete Algorithms. Society for Industrial and Applied Mathematics, USA, pp. 1027–1035.\par
Benchimol, M., Benchimol, P., Chappert, B., De La Taille, A., Laroche, F., Meunier, F., Robinet, L., 2011. Balancing the stations of a self service “bike hire” system. RAIRO-Oper. Res. 45 (1), 37–61.\par
Chemla, D., Meunier, F., Wolfler Calvo, R., 2013. Bike sharing systems: solving the static rebalancing problem. Discrete Optim. 10 (2), 120–146.\par
Dantzig, G., Fulkerson, R., Johnson, S., 1954. Solution of a large-scale traveling-salesman problem. J. Oper. Res. Soc. Am. 2 (4), 393–410.\par
Di Gaspero, L., Rendl, A., Urli, T., 2013. Constraint-based approaches for balancing bike sharing systems. In: Schulte, C. (Ed.), Principles and Practice of Constraint Programming. Springer Berlin Heidelberg, Berlin, Heidelberg, pp. 758–773.\par
Forma, I.A., Raviv, T., Tzur, M., 2015. A 3-step math heuristic for the static repositioning problem in bike-sharing systems. Transp. Res. Part B: Methodol. 71, 230–247.\par
Ho, S.C., Szeto, W.Y., 2014. Solving a static repositioning problem in bike-sharing systems using iterated tabu search. Transp. Res. Part E: Logist. Transp. Rev. 69, 180–198.\par
Ho, S.C., Szeto, W.Y., 2017. A hybrid large neighborhood search for the static multi-vehicle bike-repositioning problem. Transp. Res. Part B: Methodol. 95, 340–363.\par
Hoeffding, W., 1963. Probability inequalities for sums of bounded random variables. J. Am. Stat. Assoc. 58 (301), 13–30.\par
Hu, R., Szeto, W.Y., Ho, S.C., 2025. Repositioning in bike sharing systems with broken bikes considering on-site repairs. Transp. Res. Part E: Logist. Transp. Rev. 201, 104155.\par
Huang, D., Chen, Xinyuan, Liu, Z., Lyu, C., Wang, S., Chen, Xuewu, 2020. A static bike repositioning model in a hub-and-spoke network framework. Transp. Res. Part E: Logist. Transp. Rev. 141, 102031.\par
Liu, J., Chen, W., Sun, L., 2025. A data-driven optimization framework for static rebalancing operations in bike sharing systems. INFORMS J. Comput. 37 (5), 1369–1390.\par
Lv, C., Liu, Q., Zhang, C., Ren, Y., Zhou, H., 2024. A feature correlation reinforce clustering and evolutionary algorithm for the green bike-sharing reposition problem. Comput. Oper. Res. 166, 106627.\par
Lv, C., Zhang, C., Lian, K., Ren, Y., Meng, L., 2020. A hybrid algorithm for the static bike-sharing re-positioning problem based on an effective clustering strategy. Transp. Res. Part B: Methodol. 140, 1–21.\par
Lv, C., Zhang, C., Ren, Y., Meng, L., 2022. A fuzzy correlation based heuristic for dual-mode integrated location routing problem. Comput. Oper. Res. 146, 105923.\par
Rainer-Harbach, M., Papazek, P., Raidl, G.R., Hu, B., Kloimüllner, C., 2015. PILOT, GRASP, and VNS approaches for the static balancing of bicycle sharing systems. J. Global Optim. 63 (3), 597–629.\par
Raviv, T., Kolka, O., 2013. Optimal inventory management of a bike-sharing station. IIE Trans. 45 (10), 1077–1093.\par
Raviv, T., Tzur, M., Forma, I.A., 2013. Static repositioning in a bike-sharing system: models and solution approaches. EURO J. Transp. Logist. 2 (3), 187–229.\par
Schuijbroek, J., Hampshire, R.C., Van Hoeve, W.-J., 2017. Inventory rebalancing and vehicle routing in bike sharing systems. Eur. J. Oper. Res. 257 (3), 992–1004.\par
Shui, C.S., Szeto, W.Y., 2020. A review of bicycle-sharing service planning problems. Transp. Res. Part C: Emerging Technol. 117, 102648.\par
Szeto, W.Y., Liu, Y., Ho, S.C., 2016. Chemical reaction optimization for solving a static bike repositioning problem. Transp. Res. Part D: Transp. Environ. 47, 104–135.\par
Vidal, T., 2022. Hybrid genetic search for the CVRP: open-source implementation and SWAP* neighborhood. Comput. Oper. Res. 140, 105643.\par
Vidal, T., Crainic, T.G., Gendreau, M., Lahrichi, N., Rei, W., 2012. A hybrid genetic algorithm for multidepot and periodic vehicle routing problems. Oper. Res. 60 (3), 611–624.\par
You, P.-S., 2019. A two-phase heuristic approach to the bike repositioning problem. Appl. Math. Modell. 73, 651–667.\par
Zhang, D., Yu, C., Desai, J., Lau, H.Y.K., Srivathsan, S., 2017. A time-space network flow approach to dynamic repositioning in bicycle sharing systems. Transp. Res. Part B: Methodol. 103, 188–207.\par
\endgroup

\end{document}
